% concatenated from paper.tex
\documentclass[12pt]{article}
\usepackage{smile}

%\setlength{\oddsidemargin}{0.25in}
%\setlength{\textwidth}{6in}
%\setlength{\topmargin}{0.5in}
%\setlength{\textheight}{9in}

% === Page Layout ===
%\usepackage{abstract}
%\usepackage{authblk}
\usepackage{fullpage}
%\usepackage[margin=1in,bottom=1.2in]{geometry}
\usepackage{lscape}
\usepackage{bigints}
% === Boxes ===
\usepackage{framed}
\usepackage{mdframed}
% === Numbering ===
\usepackage{enumerate}
\usepackage[inline]{enumitem}
% === Fonts ===
%\usepackage[urw-garamond]{mathdesign}
\usepackage[T1]{fontenc}
\usepackage{moresize}
\usepackage{bm}
\usepackage{bbm}
\usepackage{dsfont}
%\usepackage{mathabx}
% === Mathematics ===
\usepackage{amsmath}
\usepackage{amssymb}
\usepackage{amsthm}
\usepackage{amsfonts}
\usepackage{stmaryrd}
\usepackage{array}
\usepackage{mathrsfs}
\usepackage{mathtools} % \prescript
%\usepackage{mathdesign}
\usepackage{extarrows}
\usepackage{stackrel}
\usepackage{relsize,exscale}
\usepackage{scalerel}
% === Paragraph Formatting ===
%\usepackage{setspace}
\usepackage[nodisplayskipstretch]{setspace}
%\usepackage{indentfirst}
% === Color ===
\usepackage{color}
\usepackage[usenames,dvipsnames]{xcolor}
% === Toys ===
\usepackage{cancel}
\usepackage{soul}
\usepackage{undertilde}
%\usepackage{ulem}
\usepackage{xfrac}
\usepackage{siunitx}
% === Figures ===
\usepackage{graphicx}
\usepackage{float}
\usepackage{rotating}
%\usepackage{subfig}
\usepackage{subcaption}
\usepackage{overpic}
\usepackage[all]{xy}
% === MetaPost ===
\DeclareGraphicsRule{*}{mps}{*}{}
% === Tikz ===
\usepackage{tikz}
\usetikzlibrary{arrows,matrix,positioning,calc,automata,patterns}
% === Tables ===
\usepackage{booktabs}
\usepackage{dcolumn}
\usepackage{multirow}
%\usepackage{slashbox}
\usepackage{diagbox}
\usepackage{tabularx}
% === Source Code ===
\usepackage{verbatim}
\usepackage{listings}
\usepackage[ruled,vlined]{algorithm2e}
%\usepackage{algorithm}
%\usepackage{algpseudocode}
\usepackage{fancyvrb}
% === Links ===
\usepackage{hyperref}
%\usepackage{url}
% === Bibliography ===
\usepackage[round]{natbib}
\usepackage{sectsty}

\hypersetup{
    bookmarks=true,         % show bookmarks bar?
    unicode=false,          % non-Latin characters in Acrobat’s bookmarks
    pdftoolbar=true,        % show Acrobat’s toolbar?
    pdfmenubar=true,        % show Acrobat’s menu?
    pdffitwindow=false,     % window fit to page when opened
    pdfstartview={FitH},    % fits the width of the page to the window
    pdftitle={My title},    % title
    pdfauthor={Author},     % author
    pdfsubject={Subject},   % subject of the document
    pdfcreator={Creator},   % creator of the document
    pdfproducer={Producer}, % producer of the document
    pdfkeywords={key1, key2}, % list of keywords
    pdfnewwindow=true,      % links in new PDF window
    colorlinks=true,        % false: boxed links; true: colored links
    linkcolor=blue,         % color of internal links (change box color with linkbordercolor)
    citecolor=blue,         % color of links to bibliography
    filecolor=blue,         % color of file links
    urlcolor=cyan           % color of external links
}

\def\scaleint#1{\vcenter{\hbox{\scaleto[3ex]{\displaystyle\int}{#1}}}}

\usepackage{stackengine}
\stackMath
\newcommand\tenq[2][1]{%
\def\useanchorwidth{T}%
\ifnum#1>1%
\stackunder[0pt]{\tenq[\numexpr#1-1\relax]{#2}}{ \scriptscriptstyle\thicksim}%
\else%
\stackunder[1pt]{#2}{ \scriptstyle\thicksim}%
\fi%
}

\makeatletter
\DeclareRobustCommand\widecheck[1]{{\mathpalette\@widecheck{#1}}}
\def\@widecheck#1#2{%
    \setbox\z@\hbox{\m@th$#1#2$}%
    \setbox\tw@\hbox{\m@th$#1%
       \widehat{%
          \vrule\@width\z@\@height\ht\z@
          \vrule\@height\z@\@width\wd\z@}$}%
    \dp\tw@-\ht\z@
    \@tempdima\ht\z@ \advance\@tempdima2\ht\tw@ \divide\@tempdima\thr@@
    \setbox\tw@\hbox{%
       \raise\@tempdima\hbox{\scalebox{1}[-1]{\lower\@tempdima\box
\tw@}}}%
    {\ooalign{\box\tw@ \cr \box\z@}}}
\makeatother

%\theoremstyle{remark}

\def\given{\,|\,}
%\def\given{\mid}
\def\biggiven{\,\big{|}\,}
\def\Biggiven{\,\Big{|}\,}
\def\tr{\mathop{\text{tr}}\kern.2ex}
\def\tZ{{\tilde Z}}
\def\tX{{\tilde X}}
\def\tY{{\tilde Y}}
\def\tR{{\tilde r}}
\def\tU{{\tilde U}}
\def\tV{{\tilde V}}
\def\tW{{\tilde W}}
\def\tN{{\tilde N}}
\def\tT{{\tilde T}}
\def\tF{{\tilde F}}
\def\P{{\mathrm P}}
\def\Q{{\mathrm Q}}
\def\U{{\mathrm U}}
\def\E{{\mathrm E}}
\def\R{{\mathbbm R}}
\def\Z{{\mathbbm Z}}
\def\N{{\mathbbm N}}
\def\C{{\mathbbm C}}
\def\cov{{\rm Cov}}
%\def\T{{ \mathrm{\scriptscriptstyle T} }}
%\def{\npreceq}{\not\preceq}
\def\d{{\mathrm d}}

\newcommand{\Zeta}{\Delta_0}

\newcommand{\sfi}{\mathsf{i}}
\newcommand{\ac}{\mathrm{ac}}
\newcommand{\sgn}{\mathrm{sgn}}
%\newcommand{\sign}{\mathrm{sign}}
\newcommand{\dCov}{\mathrm{dCov}}
\newcommand{\RdCov}{\mathrm{RdCov}}
\newcommand{\HdCov}{\mathrm{HdCov}}
\newcommand{\Pos}{\mathrm{Pos}}
\newcommand{\Neg}{\mathrm{Neg}}
\newcommand{\PosNeg}{\mathrm{PosNeg}}
\renewcommand{\Pr}{\mathrm{P}}
\newcommand{\rank}{\mathrm{rank}}
\newcommand{\dpt}{$\cdot$}
\newcommand{\textand}{,~}
\newcommand{\card}{\mathrm{card}}
\renewcommand{\vec}{\mathrm{vec}}
\newcommand{\F}{\mathrm{F}}
\newcommand{\nuc}{\star}
\newcommand{\op}{\mathrm{op}}
\newcommand{\TV}{\mathrm{TV}}
\newcommand{\KL}{\mathrm{KL}}
\newcommand{\HL}{\mathrm{HL}}
\newcommand{\Hell}{\mathrm{Hell}}
\newcommand{\Kolm}{\mathrm{Kolm}}
\newcommand{\Wass}{\mathrm{Wass}}
%\newcommand{\mDe}{\fM\Delta}
%\newcommand{\wmDe}{\widehat{\fM\Delta}}
%\DeclareMathOperator{\indep}{\rotatebox[origin=c]{90}{$\models$}}
\newcommand{\indep}{\perp    \perp}
%\newcommand{\m}{\mkern-1mu}
%\newcommand{\dm}{\mkern-2mu}
%\newcommand{\muad}{\mkern-18mu}
\newcommand{\Rot}{\mathsf{Rot}}
\newcommand{\Arc}{\mathsf{Arc}}
\newcommand{\mbinom}{\binom}
\newcommand{\zahl}[1]{\llbracket #1\rrbracket}
\newcommand\yestag{\addtocounter{equation}{1}\tag{\theequation}}
\newcolumntype{L}[1]{>{\raggedright\let\newline\\\arraybackslash\hspace{0pt}}m{#1}}
\newcolumntype{C}[1]{>{  \centering\let\newline\\\arraybackslash\hspace{0pt}}m{#1}}
\newcolumntype{R}[1]{>{ \raggedleft\let\newline\\\arraybackslash\hspace{0pt}}m{#1}}
\newcolumntype{d}[1]{D{.}{.}{#1}}
\newcolumntype{H}{>{\setbox0=\hbox\bgroup}c<{\egroup}@{}}
\newcolumntype{Z}{>{\setbox0=\hbox\bgroup}c<{\egroup}@{\hspace*{-\tabcolsep}}}
\newcolumntype{b}{X}
\newcolumntype{s}{>{\hsize=.5\hsize}X}
\def\checkmark{\tikz\fill[scale=0.4](0,.35) -- (.25,0) -- (1,.7) -- (.25,.15) -- cycle;}
%\newcommand{\cmark}{\ding{51}}
%\newcommand{\xmark}{\ding{55}}
\newcommand{\cmark}{Yes}
\newcommand{\xmark}{No}
\newcommand{\n}{^{(n)}}

\def\pms{\mspace{-1mu}{\scriptscriptstyle{\pm}}}

%\numberwithin{equation}{section}
%\theoremstyle{plain}
\newtheorem{condition}{Condition}
\newtheorem{theorem}{Theorem}
\newtheorem{lemma}{Lemma}
\newtheorem{proposition}{Proposition}[section]
\newtheorem{assumption}{Assumption}
\newtheorem{corollary}{Corollary}
\newtheorem{innercustomgeneric}{\customgenericname}
\providecommand{\customgenericname}{}
\newcommand{\newcustomtheorem}[2]{%
  \newenvironment{#1}[1]
  {%
   \renewcommand\customgenericname{#2}%
   \renewcommand\theinnercustomgeneric{##1}%
   \innercustomgeneric
  }
  {\endinnercustomgeneric}
}
\newcustomtheorem{customdefinition}{Definition}
\newcustomtheorem{customdefinitions}{Definitions}
\newcustomtheorem{customtheorem}{Theorem}
\newcustomtheorem{customassumption}{Assumption}
\newcustomtheorem{customlemma}{Lemma}
\newcustomtheorem{customexample}{Example}
\theoremstyle{definition}
\newtheorem{definition}{Definition}
\newtheorem{example}{Example}
\newtheorem{remark}{Remark}

\usepackage{enumitem}
\makeatletter
\newcommand{\mylabel}[2]{#2\def\@currentlabel{#2}\label{#1}}
\makeatother

\setcounter{tocdepth}{4}
\setcounter{secnumdepth}{4}

\graphicspath{{./fig3/}}

%\onehalfspacing
%\doublespacing
\renewcommand{\baselinestretch}{1.09}
%\setstretch{1.1}

\newcommand{\nb}[1]{{\color{red}{#1}}}
\newcommand{\hf}[1]{{\color{magenta}#1}}

\allowdisplaybreaks

\begin{document}

\setlength{\abovedisplayskip}{5pt}
\setlength{\belowdisplayskip}{5pt}
\setlength{\abovedisplayshortskip}{5pt}
\setlength{\belowdisplayshortskip}{5pt}
\hypersetup{colorlinks,breaklinks,urlcolor=blue,linkcolor=blue}



\title{\LARGE Nearest-Neighbor Radii under Dependent Sampling}

\renewcommand{\thefootnote}{\fnsymbol{footnote}}

\author{
Yuanyuan Gao$^{1,\ast}$ \quad
Yilong Hou$^{2,\ast}$ \quad
Zhexiao Lin$^{1,\ast,\dagger}$ \\[0.4em]
\small $^1$Department of Statistics, University of California, Berkeley, CA 94720, USA \\
\small $^2$Department of Biostatistics, University of California, Berkeley, CA 94720, USA
}

\date{\today}

\maketitle

\footnotetext{$^\ast$Authors are listed in alphabetical order.}
\footnotetext{$^\dagger$Corresponding author: \texttt{zhexiaolin@berkeley.edu}.}

\vspace{-1em}

\begin{abstract}
Nearest-neighbor methods are fundamental to classical and modern machine learning, yet their geometric properties are typically analyzed under independent sampling. In this paper, we study the nearest-neighbor radii under dependent sampling. We consider strong mixing dependent observations and ask whether dependence changes the scale of nearest-neighbor neighborhoods. We establish distribution-free almost sure convergence under polynomial mixing and sharp non-asymptotic moment bounds under geometric mixing. The moment bounds depend on the local intrinsic dimension rather than the ambient dimension, making the results applicable to high-dimensional data concentrated near lower-dimensional manifolds. Synthetic experiments and real-world time-series benchmarks support the theory, showing that nearest-neighbor geometry remains informative under dependence sampling.
\end{abstract}

\section{Introduction}\label{sec:intro}

Nearest-neighbor (NN) methods remain one of the most fundamental ideas in machine learning. 
They underlie classical procedures such as $k$-nearest-neighbor classification \citep{cover1967nearest} and regression \citep{stone1977consistent}, 
nonparametric density estimation and local averaging \citep{loftsgaarden1965nonparametric}, and they also appear implicitly in a much broader range of modern techniques built on local neighborhoods, including graph-based semi-supervised learning \citep{zhu2003semi}, manifold methods \citep{belkin2003laplacian}, information retrieval \citep{wang2021comprehensive}, and language modeling \citep{khandelwal2020nearest,xu2023k, xu2023nearest}. In all of these methods, a basic geometric quantity at the center of the analysis is the distance from a query point $x$ to its $k$th nearest neighbor in the sample. This random radius determines the scale of the local neighborhood around $x$, and therefore controls how much of the ambient space the method actually ``sees.''

From a statistical perspective, nearest-neighbor radii determine the central bias--variance tradeoff of local methods \citep{stone1977consistent, biau2015lectures}. If the radius is too large, a local estimator averages over points that are geometrically far from the query and thus introduces bias. If the radius is too small, the estimator uses too few effective observations and becomes unstable. In density estimation, the $k$-NN radius is the random bandwidth that determines the volume of the adaptive neighborhood. In regression and classification, it controls the locality of smoothing. In manifold and graph-based methods, it determines neighborhood connectivity and the scale at which local geometry is resolved. Understanding the scaling of the $k$-NN radius is therefore a prerequisite for understanding the performance of a wide class of machine learning algorithms.

In the classical i.i.d.\ setting, the behavior of $k$-NN radius is well understood: under standard local mass conditions, the canonical scale is $(k/n)^{1/d}$, or more generally, $(k/n)^{1/s}$ when the local geometry is determined by an intrinsic dimension $s$ \citep{loftsgaarden1965nonparametric,biau2015lectures,kpotufe2011k}. However, many learning problems do not arise from independent samples. Data are often collected sequentially, spatially, or through interaction with a dynamical environment. Examples include time-series forecasting, reinforcement learning from trajectories, sequential recommendation, event streams, dependent tabular logs, and spatiotemporal sensing \citep{mohri2010stability,tu2024learning}. In such settings, observations can be strongly dependent even when their marginal distribution is stable. This raises a basic question:

\begin{center}
\emph{How does dependence affect the geometry of nearest-neighbor neighborhoods?}
\end{center}

This question is both mathematically natural and practically important. A large body of machine learning theory relies on i.i.d.\ nearest-neighbor geometry either explicitly or implicitly. Yet once dependence is present, standard concentration arguments no longer apply directly, and even basic geometric quantities such as the $k$-NN radius become more delicate. If dependence were to fundamentally change the scaling of nearest-neighbor radii, then the geometry underlying many local learning procedures would also change. On the other hand, if the canonical radius scale remains unchanged under weak dependence, this would help justify the use of nearest-neighbor and local methods beyond the i.i.d. regime.

In this paper, we study the nearest-neighbor radii under dependent sampling. We consider a triangular array of $\R^d$-valued observations with a common marginal distribution and strong mixing dependence along each row. This framework is general enough to cover the usual stationary-sequence setting as a special case, while also allowing the dependence structure to vary with $n$. Our object of study is the radius $R_{n,k}(x) := \| \mathcal X_{(k)}^n(x)-x\|$,
that is, the distance from a fixed query point $x$ to its $k$th nearest neighbor among the $n$ observations in the row.

Our first result is almost sure convergence: under polynomial $\alpha$-mixing, we show that the nearest-neighbor radius converges almost surely to the distance from $x$ to the support of the marginal law. This gives a dependent-sampling analogue of the classical consistency result and shows that weak dependence does not change the limiting geometry of nearest-neighbor radii. We then develop a non-asymptotic theory under geometric mixing. Under a local lower-mass condition of order $s$, we establish a tail bound for $R_{n,k}(x)$ and derive the sharp moment upper bound $\E[R_{n,k}(x)^p] \lesssim (k/n)^{p/s}$ in the standard regime $k \gtrsim \log n$. Under a matching local upper-mass condition of order $s$, we also establish the matching lower bound and hence obtain the exact order $\E[R_{n,k}(x)^p] \asymp (k/n)^{p/s}$. Thus, geometric dependence changes the concentration mechanism, but not the fundamental exponent.

A key feature of our theory is that it is formulated in terms of a local exponent $s>0$, rather than the ambient dimension $d$. This point is especially relevant for modern machine learning, where high-dimensional data often concentrate near lower-dimensional latent structures \citep{fefferman2016testing,ansuini2019intrinsic}. Our results therefore support the common regime in which the effective geometry around the query point is lower-dimensional than the ambient space. In this sense, our results are not only about dependence, but also about how dependence interacts with local geometric complexity.

\paragraph*{Our contributions.}
The main contributions are as follows.
\begin{enumerate}
    \item We establish a \emph{distribution-free almost sure convergence theorem} for the $k$-NN radius under polynomial $\alpha$-mixing. This extends the classical i.i.d. consistency result to dependent triangular arrays.

    \item Under geometric $\alpha$-mixing and a local lower-mass condition of order $s$, we establish a \emph{non-asymptotic tail bound} for the $k$-NN radius and derive the \emph{sharp moment upper bound} $\E[R_{n,k}(x)^p] \lesssim (k/n)^{p/s}$ for all $p>0$ in the standard dependent regime $k \gtrsim \log n$.

    \item Under a matching local upper-mass condition, we establish a \emph{matching lower bound} $\E[R_{n,k}(x)^p] \gtrsim (k/n)^{p/s}$, and hence obtain the \emph{exact-order characterization} $\E[R_{n,k}(x)^p] \asymp (k/n)^{p/s}$.

    \item Our theory is stated in terms of a \emph{local geometric exponent} $s$, which yields an intrinsic-dimension interpretation and is relevant for modern high-dimensional learning problems.
\end{enumerate}

We believe the main message is simple and useful: under weak dependence, nearest-neighbor geometry remains essentially the same as in the independent case at the level of its canonical scale. This provides theoretical support for transferring geometric intuition and local-method design principles from the i.i.d.\ world to substantially broader dependent settings.

\subsection{Related work}
\label{sec:related}

\paragraph*{Applications of nearest-neighbor methods.} In recent years, nearest-neighbor methods have re-emerged in modern machine learning as a flexible mechanism for retrieval, adaptation, and locality in learned representation spaces. Examples include self-supervised representation learning \citep{avdiukhin2024embedding,long2025mnn}, tabular deep learning \citep{gorishniy2023tabr}, vector search and approximate nearest-neighbor infrastructure \citep{chen2021spann,yang2024cspg,jaasaari2024lorann,wang2021comprehensive}, and robust retrieval and vision \citep{wu2022retrievalguard,nizan2024k}. Beyond core machine learning, nearest-neighbor structure also plays an important role in causal inference through matching estimators \citep{abadie2006large,abadie2011bias,lin2023estimation,lin2025regression,cattaneo2025rosenbaum}, in rank-based and graph-based measures of dependence \citep{chatterjee2021new,azadkia2021simple,han2021extensions,lin2022limit,lin2023boosting,lin2024failure, shi2024azadkia, han2024azadkia}, in information-theoretic functional estimation via $k$-NN entropy and mutual-information estimators \citep{kozachenko1987sample,kraskov2004estimating}, and in multivariate nonparametric testing based on geometric graphs \citep{friedman1979multivariate,henze1999multivariate}. This breadth of applications makes the geometric behavior of nearest-neighbor neighborhoods a foundational issue.

\paragraph*{Theory of nearest-neighbor methods under dependence.} Despite the wide range of applications, previous literature on the theory under dependence has mostly analyzed estimators built from nearest-neighbor neighborhoods, rather than the geometry of the neighborhoods themselves. Examples include nearest-neighbor classification with dependent training sequences \citep{holst2001nearest} and rates for nearest-neighbor estimation under more general non-i.i.d.\ sampling schemes \citep{kulkarni1995rates}. More broadly, the literature on learning from dependent observations has developed both positive and negative results for nonparametric prediction under mixing, ergodic, and other non-i.i.d.\ assumptions \citep{yu1994rates, nobel1999limits,mohri2010stability}. By contrast, our focus is on the radius itself as a primary object. This shift in perspective is useful because the nearest-neighbor radius serves as the random local scale in many downstream methods, so controlling it directly provides a geometric ingredient for analyzing local learning methods.


\section{Setup}
\label{sec:setup}

We study the geometry of nearest neighbors under dependent sampling. Let $[X_{n,i}: 1 \le i \le n,\ n \ge 1]$ be a triangular array of $\R^d$-valued random variables defined on a common probability space $(\Omega,\cF,\P)$, where the ambient dimension $d \ge 1$ is fixed. For each row $n$, the observations $X_{n,1},\dots,X_{n,n}$ could be independent, but we assume that they share a common marginal law $\mu$ on $\R^d$. Thus, for every $n \ge 1$ and every $i \in \{1,\dots,n\}$, $X_{n,i} \sim \mu$.

This formulation distinguishes two roles in our analysis. The common marginal distribution $\mu$ determines the local geometry of the marginal distribution, while the dependence structure controls how much effective information is contained in a sample of size $n$. Allowing the dependence structure to vary with $n$ leads naturally to a triangular-array framework, which is flexible enough to cover the usual stationary-sequence setting as a special case and at the same time convenient for stating non-asymptotic results uniformly in $n$.

\paragraph*{Triangular-array strong mixing.} For integers $1 \le k \le m \le n$, let ${}^n\cF_k^m := \sigma(X_{n,i} : k \le i \le m)$ denote the $\sigma$-field generated by the segment $(X_{n,k},\dots,X_{n,m})$ in the $n$th row. Following the standard strong-mixing framework, define the mixing coefficients of the triangular array by
\[
\alpha_m
:=
\sup_{n \ge 1}\ \sup_{1 \le k \le n-m}\ 
\sup\Big\{
\big| \P(A \cap B) - \P(A)\P(B) \big|
:\ 
A \in {}^n\cF_1^k,\ 
B \in {}^n\cF_{k+m}^n
\Big\},
\qquad m \ge 1.
\]
The array is said to be \emph{strong mixing} (or \emph{$\alpha$-mixing}) if
\[
\alpha_m \to 0
\qquad\text{as } m \to \infty.
\]

The strong-mixing framework is classical in probability and time-series analysis \citep{bradley2005basic}. This definition requires a uniform decay of dependence across all rows of the array. Intuitively, if two collections of variables are separated by at least $m$ indices within the same row, then their dependence becomes weak when $m$ is large. This is the dependence notion used throughout the paper.

The triangular-array formulation is more general than working with a single stationary process $[X_i]_{i\ge1}$. It allows the dependence structure to vary with $n$, while keeping the marginal distribution fixed. We state finite-sample results uniformly in $n$, which also covers the usual time-series setting as a special case by taking $X_{n,i} = X_i$ for $1 \le i \le n$. Therefore, readers who prefer the standard sequential viewpoint may regard our framework as a generalization.

\paragraph*{Nearest-neighbor radii.}
Fix a query point $x \in \R^d$. For each $n \ge 1$ and $k \in \{1,\dots,n\}$, let
\[
\mathcal X_{(k)}^n(x) \in \{X_{n,1},\dots,X_{n,n}\}
\]
denote the $k$th nearest neighbor of $x$ among the sample points in the $n$th row, with ties broken arbitrarily. We introduce the $k$-NN radius:
\[
  R_{n,k}(x) := \|\mathcal X_{(k)}^n(x)-x\|,
\]
together with the local counting process:
\[
  N_n(x,r)
  :=
  \sum_{i=1}^n \mathbf{1}\{\|X_{n,i}-x\| \le r\}.
\]
These two quantities are equivalent through the relationship
\[
R_{n,k}(x) > r
\quad\Longleftrightarrow\quad
N_n(x,r) < k.
\]
This observation is the starting point for our entire analysis: upper bounds on the radius are equivalent to lower bounds for the number of sample points falling in $B(x,r)$. Under independence, $N_n(x,r)$ is binomial with parameters $n$ and $\mu(B(x,r))$, so its behavior can be analyzed using classical binomial tail bounds. Under dependence, however, this analysis is not applicable since the indicators $\mathbf{1}\{\| X_{n,i}-x\| \le r\}$ are dependent, and hence the local counting process is not binomial. Controlling $R_{n,k}(x)$ therefore requires concentration arguments for dependent sequences, which is one of the main technical challenges.

\paragraph*{Support and target limit.}
The asymptotic behavior of $R_{n,k}(x)$ depends on the position of the query point relative to the support of $\mu$. Define
\[
{\rm supp}(\mu)
:=
\{x \in \R^d : \mu(B(x,r)) > 0 \text{ for all } r>0\},
\]
where $B(x,r)$ denotes the closed Euclidean ball centered at $x$ with radius $r$. We also define the distance from $x$ to the support of $\mu$ by
\[
r_x := \inf_{y \in {\rm supp}(\mu)} \| y-x \|.
\]
When $x \in {\rm supp}(\mu)$, one has $r_x=0$; otherwise $r_x$ is the deterministic limit that the $k$-NN radius converges to under mild assumptions.

% The first part of the paper establishes almost sure convergence of $R_{n,k}(x)$ to $r_x$ under polynomial strong mixing. The sharper finite-sample theory requires geometric mixing together with local mass conditions on $\mu$ around the query point $x$. Under these conditions, we study both the tail probability $\P(R_{n,k}(x) > r)$ and the moments $\E[R_{n,k}(x)^p]$ for $p>0$. Our goal is to characterize how these quantities scale with $(k,n)$ and with the local geometry of $\mu$ near $x$.

% \paragraph*{Notation.}
% For any integers $n,d\ge 1$, let $\zahl{n}:= \{1,2,\ldots,n\}$, $n!$ be the factorial of $n$, and $\R^d$ be the $d$-dimensional real space. A set consisting of distinct elements $x_1,\dots,x_n$ is written as either $\{x_1,\dots,x_n\}$ or $\{x_i\}_{i=1}^{n}$. The corresponding sequence is denoted by $[x_1,\dots,x_n]$ or $[x_i]_{i=1}^{n}$. Let $\ind(\cdot)$ denote the indicator function.
% For any $a,b \in \R$, write $a \vee b = \max\{a,b\}$ and $a \wedge b = \min\{a,b\}$.
% %For any two real sequences $\{a_n\}$ and $\{b_n\}$, write $a_n \lesssim b_n$ (or equivalently, $b_n \gtrsim a_n$ or $a_n=O(b_n)$) if there exists a universal constant $C>0$ such that $a_n/b_n \le C$ for all sufficiently large $n$, and write $a_n \prec b_n$ (or equivalently, $b_n \succ a_n$ and $a_n=o(b_n)$) if $a_n/b_n \to 0$ as $n$ goes to infinity. We write $a_n\asymp b_n$ if both $a_n\lesssim b_n$ and $b_n\lesssim a_n$ holds.
% For any two real sequences $\{a_n\}$and $\{b_n\}$, write $a_n \lesssim b_n$ (or equivalently, $b_n \gtrsim a_n$) if there exists a universal constant $C>0$ such that $a_n/b_n \le C$ for all sufficiently large $n$. We write $a_n = O(b_n)$ if $\lvert a_n \rvert \lesssim b_n$. We use  $\stackrel{\sf a.s.}{\longrightarrow}$ to denote convergence almost surely. Denote the closed ball in $\R^d$ centered at $x$ with radius $\delta$ by $B_{x,\delta}$.

\section{Main results}\label{sec:main}

We now present the main results for the $k$-NN radius under dependent sampling. As discussed in Section~\ref{sec:setup}, the behavior of $R_{n,k}(x)=\| \mathcal X_{(k)}^n(x)-x\|$ depends on both the dependence structure of the sample and the local geometry of the marginal law around the query point $x$. The almost sure limit is determined by the position of $x$ relative to the support of $\mu$, whereas the convergence rate depends on how the mass of $\mu$ accumulates near $x$. 

% Recall that the support of $\mu$ is denoted by ${\rm supp}(\mu)$, and that $r_x$ denotes the distance from $x$ to ${\rm supp}(\mu)$. In particular, if $x \in {\rm supp}(\mu)$, then $r_x=0$; otherwise, $r_x$ is the deterministic limit to which the $k$-NN radius converges.

\subsection{Distribution-free almost sure convergence}

Our first theorem shows that the almost sure convergence of the $k$-NN radius to $r_x$ holds under weak dependence.

\begin{theorem}\label{thm:as}
Assume that the triangular array $[X_{n,i}]_{1 \le i \le n,\ n \ge 1}$ is $\alpha$-mixing with coefficients $[\alpha_m]_{m \ge 1}$ satisfying $\alpha_m \lesssim m^{-\gamma}$ for some $\gamma>1$. If $k/n \to 0$, then
\[
\|\mathcal X_{(k)}^n(x)-x\| \stackrel{\sf a.s.}{\longrightarrow} r_x.
\]
In particular, if $x \in {\rm supp}(\mu)$ and $k/n \to 0$, then
\[
\|\mathcal X_{(k)}^n(x)-x\| \stackrel{\sf a.s.}{\longrightarrow} 0.
\]
\end{theorem}

Theorem~\ref{thm:as} may be viewed as the dependent-sampling analogue of the classical i.i.d. consistency result \citep[Lemma 2.2]{biau2015lectures}. It shows that weak dependence, quantified here through a polynomial mixing condition, does not change the limiting geometric behavior of nearest-neighbor radii. Importantly, the result remains distribution-free: no density, smoothness, or local mass condition is required. The only assumption is that dependence decays polynomially fast along each row of the triangular array.


\subsection{From almost sure convergence to moment bounds}

We next turn to non-asymptotic moment results. For these stronger results, one must quantify how much probability mass lies near the query point $x$.

Rather than working with the ambient dimension $d$, we state the results using a general local exponent $s>0$. This formulation is more flexible and captures the possibility that the effective local dimension near $x$ is smaller than the ambient dimension. The standard Euclidean case is recovered by setting $s=d$.

For the non-asymptotic results, we impose the stronger assumption that the mixing coefficients decay exponentially fast.
\begin{assumption}[Geometric $\alpha$-mixing]\label{ass:geom}
There exists $\gamma>0$ such that $\alpha_m \le e^{-\gamma m}$ for all $m \ge 1$.
\end{assumption}
  
Assumption~\ref{ass:geom} strengthens the assumption in Theorem~\ref{thm:as} from polynomial to exponential decorrelation. This stronger assumption enables non-asymptotic concentration via blocking arguments. The parameter $\gamma$ quantifies the strength of dependence: larger $\gamma$ corresponds to faster decorrelation, and in the limit of vanishing dependence, one recovers behavior close to the i.i.d. case.

Geometric mixing is a standard assumption in time-series theory and is satisfied by many classical models under mild regularity conditions \citep{doukhan2012mixing}. In particular, stationary ARMA processes are known to be strongly mixing with exponential rate under standard stability assumptions \citep{mokkadem1988mixing}, which makes geometric mixing a natural benchmark for studying dependent nearest-neighbor geometry. 
  
\begin{assumption}[Local lower mass of order $s$ at $x$]\label{ass:lower}
There exist constants $c_->0$, $r_0>0$, and $s>0$ such that $\mu(B(x,r)) \ge c_- r^s$ for all $0<r\le r_0$.
\end{assumption}
  
Assumption~\ref{ass:lower} guarantees that the ball $B(x,r)$ contains enough probability mass for its radius-$r$ neighborhood to capture roughly $nr^s$ sample points. This lower bound is the key ingredient for the upper bound of $R_{n,k}(x)$.
  

\begin{assumption}[Local upper mass of order $s$ at $x$]\label{ass:upper}
There exist constants $c_+>0$, $r_0>0$, and $s>0$ such that $\mu(B(x,r)) \le c_+ r^s$ for all $0<r\le r_0$.
\end{assumption}

Assumption~\ref{ass:upper} matches Assumption~\ref{ass:lower} and prevents the distribution from placing too much mass near $x$, and is needed to establish a matching lower bound on the radius.
  
\begin{assumption}[Compact support]\label{ass:compact}
The support of $\mu$ is contained in a compact set of diameter at most $D$.
\end{assumption}

Assumption~\ref{ass:compact} is for technical simplicity rather than a geometric limitation. The local behavior of the $k$-NN radius is determined by Assumptions~\ref{ass:lower} and~\ref{ass:upper}, which quantify how the mass of $\mu$ accumulates near the query point $x$. By contrast, compact support is used to control the far tail of $R_{n,k}(x)$ when converting non-asymptotic tail bounds into moment bounds. Compact support ensures that the radius is uniformly bounded, which keeps the moment argument simpler. One could replace the compact support by a tail or moment condition on $\mu$ at the cost of a more technical proof.

\paragraph*{Interpretation of the local exponent $s$.}
The exponent $s$ is the \emph{local geometric dimension} around the query point $x$. Assumptions~\ref{ass:lower} and~\ref{ass:upper} say that, for sufficiently small radii, $\mu(B(x,r)) \asymp r^s$. Thus, the amount of probability mass captured by a small neighborhood of $x$ grows like the volume of an $s$-dimensional ball. In the standard Euclidean case with a positive density in $\R^d$, one has $s=d$. More generally, $s$ may be strictly smaller than $d$, e.g., when the distribution is concentrated near a lower-dimensional manifold, subspace, or other latent geometric structure. Our theory is therefore local and intrinsic: the scaling of the $k$-NN radius is determined by the intrinsic dimension of the distribution near $x$, rather than by the ambient Euclidean dimension. 

Throughout the sharp upper-bound results, we work in the regime $K_0 \log n \le k \le \kappa n$, for some fixed $K_0>0$ and $0 < \kappa < c_- r_0^s/8$. The upper bound ensures that the radius remains within the local regime $r\le r_0$, where Assumption~\ref{ass:lower} applies. The lower bound $k \gtrsim \log n$ is due to dependence: it arises because the proof decouples the sample into blocks of length $\asymp \log n$, as is standard in Bernstein-type concentration arguments for geometric mixing sequences \citep{liebscher1996strong,merlevede2009bernstein}. In contrast, the lower bound on the radius does not require $k \gtrsim \log n$. 

\paragraph*{On the regime $k \gtrsim \log n$.} The lower bound $k \gtrsim \log n$ is the natural regime in which our blocking-based concentration argument resolves the dependence structure. It reflects the cost of decoupling geometric mixing observations into approximately independent blocks, rather than a change in the underlying nearest-neighbor geometry. In particular, our results suggest that dependence affects concentration through a logarithmic penalty, while the canonical radius scale remains $(k/n)^{1/s}$.

\subsection{Non-asymptotic upper bound}

We first establish a tail bound on $R_{n,k}(x)$.
\begin{theorem}[Non-asymptotic tail bound]\label{thm:tail}
Assume Assumptions~\ref{ass:geom}, \ref{ass:lower}, and \ref{ass:compact}. Define
\[
  b_n := \Big\lceil \frac{8}{\gamma}\log n \Big\rceil,
  \qquad
  r_*(n,k) := \Big(\frac{8k}{c_- n}\Big)^{1/s}.
\]
Then there exists a constant $c_0>0$, depending only on $\gamma$, such that for every integer $j \ge 0$ satisfying $2^j r_*(n,k) \le r_0$, we have
\[
  \P\big(R_{n,k}(x) > 2^j r_*(n,k)\big)
  \le
  4\exp \Big(-c_0\,2^{js}\,\frac{k}{\log n}\Big)
  + C n^{-7},
\]
for all sufficiently large $n$.
\end{theorem}

Theorem~\ref{thm:tail} shows that the $k$-NN radius is concentrated around the canonical geometric scale $(k/n)^{1/s}$. In particular, Theorem~\ref{thm:tail} separates the roles of geometry and dependence: the exponent $1/s$ is determined by the local mass behavior of $\mu$ near $x$, while the dependence structure enters only through the additional $\log n$ factor in the concentration exponent, which is the price of the blocking argument under geometric mixing. Thus, although dependence weakens concentration relative to the i.i.d.\ case, it does not change the underlying nearest-neighbor scale. Theorem~\ref{thm:tail} implies a high-probability upper bound:
\[
  R_{n,k}(x)
  \lesssim
  \Big(\frac{k+t\log n}{n}\Big)^{1/s}
\]
with probability at least $1-e^{-ct}$, for some constant $c>0$.

Integrating the tail bound yields the sharp moment upper bound.


\begin{theorem}[Sharp moment upper bound]\label{thm:moment}
Assume Assumptions~\ref{ass:geom}, \ref{ass:lower}, and \ref{ass:compact}. Let $p>0$. Then there exists a constant $C_{p,s}>0$ such that, whenever $K_0\log n \le k \le \kappa n$, we have
\[
  \E\big[R_{n,k}(x)^p\big]
  \le
  C_{p,s}\Big(\frac{k}{n}\Big)^{p/s}.
\]
\end{theorem}

If $\mu$ admits a density, then by the Lebesgue differentiation theorem, Assumption~\ref{ass:lower} holds with $s=d$ for $\mu$-almost every interior point at which the density is positive. Consequently, Theorem~\ref{thm:moment} applies for $\mu$-almost all such $x$. Theorem~\ref{thm:moment} is applicable in downstream analyses, since many local learning bounds depend on moments of the random neighborhood. We discuss this further in the discussion section.

\subsection{Matching non-asymptotic lower bound}

To show that the exponent $p/s$ is optimal, we next establish a matching lower bound. Unlike the upper bound, the lower bound is purely local and does not require any dependence condition.

\begin{theorem}[Matching lower bound]\label{thm:lower}
Assume Assumption~\ref{ass:upper}. Let $p>0$. Then whenever $(\frac{k}{2c_+ n})^{1/s} \le r_0$, we have
\[
  \E\big[R_{n,k}(x)^p\big]
  \ge
  \frac12\Big(\frac{k}{2c_+ n}\Big)^{p/s}.
\]
\end{theorem}

Theorem~\ref{thm:lower} is conceptually simple but important: if a ball of radius $r$ contains at most order $r^s$ probability mass, then one cannot expect the $k$th nearest neighbor to lie much closer than $(k/n)^{1/s}$. Combining the upper and lower bounds yields the following exact-order characterization.
\begin{theorem}[Exact-order characterization]\label{thm:exact}
Assume Assumptions~\ref{ass:geom}, \ref{ass:lower}, \ref{ass:upper}, and \ref{ass:compact}, with the same exponent $s>0$. Then for every $p>0$ there exist constants $0 < c_{p,s} \le C_{p,s} < \infty$ such that, for all sufficiently large $n$ and all $K_0\log n \le k \le \kappa n$, we have
\[
  c_{p,s}\Big(\frac{k}{n}\Big)^{p/s}
  \le
  \E\big[R_{n,k}(x)^p\big]
  \le
  C_{p,s}\Big(\frac{k}{n}\Big)^{p/s}.
\]
\end{theorem}

Theorem~\ref{thm:exact} shows that the canonical nearest-neighbor radii scale remains unchanged under geometric mixing. The dependence changes constants and concentration behavior, but not the fundamental exponent.


\section{Experiments}

\subsection{Moment convergence rate} \label{exp1:Moment convergence rate}

We test whether the moment-rate upper bound suggested by the proved $\alpha$-mixing setting remains visible under broader dependence structures. The implementation compares the i.i.d.\ benchmark with three dependent sequence models: a linear Gaussian state-space model, a hidden Markov model, and a stationary Gaussian-process sequence with a radial basis function (RBF)-like kernel generated by fast Fourier transform (FFT). All observations are marginally transformed to the cube $[0,1]^d$.

We sample 1000 points uniformly from the interior of the cube and report the mean $p$-th moment of the $k$-NN radius. For each $d\in\{1,3,5\}$, we choose $p$ such that $p/d\in\{1,2,3,4,5\}$. We use exponentially growing, dimension-dependent sample sizes $n(d)=100\cdot 2^d\cdot m$ for $m\in\{1,2,4,8,16,32\}$. The number of NNs is chosen as $k=\lfloor n^\beta\rceil$ with $\beta$ on a 20-point grid in $[0.1,0.9]$ for slope fitting. The slope is the coefficient $b$ from $\log M = a + b\log(k/n)$, with theoretical slope $p/d$ for the i.i.d.\ and $\alpha$-mixing setting.

\begin{figure}[H]
    \centering
    \includegraphics[width=\textwidth]{uniform_interior_mean_slope_summary_grid.pdf}
    \caption{Slope summary for Experiment~1: rows correspond to $d=1,3,5$, columns correspond to the three dependent model classes, and the shared legend compares the theoretical line $p/d$, the i.i.d.\ benchmark, and weak/medium/strong dependence.}
    \label{fig:exp1-slope-grid}
\end{figure}

Across the three data- generating processes (linear state-space, hidden Markov, and FFT Gaussian process), we impose weak/medium/strong dependence via model-specific persistence parameters. Figure~\ref{fig:exp1-slope-grid} shows the expected pattern: weak dependence stays close to the i.i.d.\ baseline and the theoretical rate $p/d$, while medium and strong dependence diverge.

\subsection{Entropy estimation} \label{exp2: Entropy estimation}

This experiment examines the intrinsic-dimension message of the nearest-neighbor theory in a representation-learning setting. We state our results in terms of a local intrinsic dimension $s$. A Kozachenko--Leonenko entropy estimate is built from the NN radii and the corresponding $s$-dimensional volume correction. Thus, if a learned representation preserves the local $s$-dimensional geometry, the entropy estimator should converge to the entropy of the intrinsic process rather than behave as though the ambient dimension is a larger $d$. We consider the setup in which observations have ambient dimension $d$, but the true data-generating process has intrinsic dimension $s$.

The experiment generates an $s$-dimensional Gaussian AR(1) latent process, whose entropy is available in closed form as $H_s=(s/2)\log(2\pi e)$, and maps it into $20$-dimensional observed data through a linear augmentation. Figure~\ref{fig:exp2-entropy} reports the results for $s\in\{1,3,5\}$ under three temporal-dependence regimes for the sample sequence: i.i.d. ($\rho=0$), weak AR(1) dependence ($\rho=0.3$), and strong AR(1) dependence ($\rho=0.6$), using the current run configuration with $1000$ replications and number of NNs $k=\max\{2,\lceil n^{0.1}\rceil\}$. The oracle estimator on the latent $s$-dimensional data and the estimator based on the top $s$ coordinates of a $5$-component principal component analysis (PCA) representation both move toward the true entropy as $n$ grows. Our reason for expecting the entropy estimator based on PCA to match the intrinsic-dimension behavior is that, in the implementation, the latent coordinates are independent across coordinates even though each coordinate follows the same AR(1) latent process. The $20$-dimensional augmented observations are highly correlated in coordinates. PCA projects the observation to a low-dimensional representation, so the leading PCA coordinates are much closer to the intrinsic geometry and preserve the same weak-dependence structure. The comparison therefore supports the conclusion that PCA recovers the intrinsic entropy behavior and that weak dependence mainly changes constants rather than the qualitative convergence pattern; implementation details and numerical tables are left to the Appendix.

\begin{figure}[H]
    \centering
    \includegraphics[width=1\textwidth]{experiment2_entropy_schedule_comparison.pdf}
    \caption{Experiment~2 entropy estimation. The dotted horizontal line is the true entropy. Colors encode the dependence level and markers distinguish the oracle intrinsic estimate from the PCA estimate.}
    \label{fig:exp2-entropy}
\end{figure}



\subsection{Real-data evaluation}\label{exp3: Real-data_evaluation}

% In this section, we complement the synthetic experiments with real-world time-series benchmarks. We compare two baselines: a raw $k$-nearest-neighbors (Raw-$k$NN) method operating directly in standardized observation space, and a PCA $k$-nearest-neighbors (PCA-$k$NN) method that first projects each input onto a lower-dimensional representation using PCA before applying nearest-neighbor prediction. We evaluate Raw-$k$NN and PCA-$k$NN on three real-data tasks: long-horizon forecasting, short-horizon forecasting, and classification. The datasets are from the Time Series Pile \citep{goswami2024moment}, a large public collection of time-series datasets compiled from multiple benchmark repositories, including the Informer long-horizon forecasting datasets \citep{zhou2021informer}, the Monash Time Series Forecasting Archive \citep{godahewa2021monash}, and the UCR Time Series Archive \citep{dau2019ucr}. We follow the benchmark setup and compare our results against MOMENT, a family of pre-trained transformer-based time-series foundation models for general-purpose time-series analysis \citep{goswami2024moment}. In both models, the input features are standardized before nearest-neighbor prediction. We tune the number of neighbors $k \in \{1,3,5,7,9,15,21,31\}$, the neighbor weighting scheme (uniform/distance), and the distance metric (Manhattan/Euclidean). For PCA-$k$NN, we additionally tune the embedding dimension over $\{8,16,32,48,64\}$.

In this section, we complement the synthetic experiments with real-world time-series benchmarks. We compare two baselines: a raw $k$-nearest-neighbors (Raw-$k$NN) method operating directly in standardized observation space, and a PCA $k$-nearest-neighbors (PCA-$k$NN) method that first projects each input onto a lower-dimensional representation using PCA before applying nearest-neighbor prediction. We evaluate Raw-$k$NN and PCA-$k$NN on three real-data tasks: long-horizon forecasting, short-horizon forecasting, and classification. The datasets are from the Time Series Pile \citep{goswami2024moment}, a large public time-series benchmark collection that includes the Informer long-horizon forecasting datasets \citep{zhou2021informer}, the Monash Time Series Forecasting Archive \citep{godahewa2021monash}, and the UCR Time Series Archive \citep{dau2019ucr}. We follow the benchmark setup and compare our results against MOMENT, a family of pre-trained time-series foundation models \citep{goswami2024moment}. In both models, the input features are standardized before nearest-neighbor prediction. We tune the number of neighbors $k \in \{1,3,5,7,9,15,21,31\}$, the neighbor weighting scheme (uniform/distance), and the distance metric (Manhattan/Euclidean). For PCA-$k$NN, we additionally tune the embedding dimension over $\{8,16,32,48,64\}$.

Table~\ref{tab:real_data_results} summarizes average performance across datasets for long-horizon forecasting, short-horizon forecasting, and classification. Figure~\ref{fig:long_horizon_mse} reports long-horizon forecasting mean squared error (MSE) separately for each dataset and forecast horizon. Figure~\ref{fig:short_horizon_vs_classification} provides complementary views of short-horizon forecasting and classification: its left panel shows short-horizon forecasting symmetric mean absolute percentage error (sMAPE) under both the within-group and source-to-target transfer settings, while its right panel shows the distribution of classification accuracy with MOMENT summary statistics included as reference lines. Together, these results are supportive in two main ways. First, Raw-$k$NN and PCA-$k$NN remain informative across all three task families. In long-horizon forecasting, the two methods achieve average MSE $0.628$ and $0.627$, respectively, compared with $0.617$ for MOMENT. In short-horizon forecasting, both methods produce nontrivial predictions in both the within-group and source-to-target transfer settings, although errors increase substantially in the transfer setting. In classification, both methods achieve mean accuracy $0.748$, compared with $0.794$ for MOMENT. These findings indicate that nearest-neighbor structure remains informative on real dependent data. Second, Raw-$k$NN and PCA-$k$NN achieve similar performance across tasks, consistent with adaptation to the intrinsic dimension of dependent data even under dependence. Their relative advantage varies across datasets, but neither method consistently dominates the other. 

Although MOMENT generally outperforms the Raw-$k$NN and PCA-$k$NN baselines, this does not contradict the theory. Both nearest-neighbor methods remain effective on real dependent data, and PCA-$k$NN performs comparably to Raw-$k$NN even after reducing the input dimension. This shows that dimension reduction can preserve the local structure needed for neighborhood-based prediction.





% \begin{table}[H]
% \centering
% \caption{Comparison of Raw-$k$NN, PCA-$k$NN, and MOMENT on real-data tasks under a common benchmark setup. Performances are averaged over datasets within each task. MOMENT results are taken from \cite{goswami2024moment}.}
% \label{tab:real_data_results}
% {\footnotesize
% \begin{tabular}{C{6.5cm}|C{1.5cm}|C{2cm}|C{2cm}|C{2cm}}
% \hline
% \textbf{Task} & \textbf{Metric} & \textbf{Raw-$k$NN} & \textbf{PCA-$k$NN} & \textbf{MOMENT} \\
% \hline
% Classification
% & Accuracy & 0.748 & 0.748 & 0.794 \\
% % \hline
% Long-horizon forecasting
% & MSE & 0.628 & 0.627 & 0.617 \\
% % \hline
% Short-horizon forecasting (transfer)
% & \multirow{2}{*}{sMAPE} & 27.210 & 27.368 & 14.525 \\
% Short-horizon forecasting (within-group)
% &  & 16.200 & 16.488 & 14.525 \\
% \hline
% \end{tabular}
% }
% \end{table}

% \begin{table}
%   \caption{Comparison of Raw-$k$NN, PCA-$k$NN, and MOMENT on real-data tasks under a common benchmark setup. Performance is averaged over datasets within each task; MOMENT results are taken from \cite{goswami2024moment}.}
%   \label{tab:real_data_results}
%   \centering
%   \footnotesize
%   \begin{tabular}{@{}lcccc@{}}
%     \toprule
%     Task & Metric & Raw-$k$NN & PCA-$k$NN & MOMENT \\
%     \midrule
%     Long-horizon forecasting & MSE & 0.628 & 0.627 & 0.617 \\
%     Classification & Accuracy & 0.748 & 0.748 & 0.794 \\
%     Short-horizon forecasting (within-group) & sMAPE & 16.200 & 16.488 & 14.525 \\
%     Short-horizon forecasting (transfer) & sMAPE & 27.210 & 27.368 & 14.525 \\
%     \bottomrule
%   \end{tabular}
% \end{table}

\begin{table}
  \caption{Comparison of Raw-$k$NN, PCA-$k$NN, and MOMENT on real-data tasks under a common benchmark setup. Performance is averaged over datasets within each task; MOMENT results are taken from \cite{goswami2024moment}.}
  \label{tab:real_data_results}
  \centering
  \footnotesize
  \begin{tabular}{@{}lcccc@{}}
    \toprule
    Task & Metric & Raw-$k$NN & PCA-$k$NN & MOMENT \\
    \midrule
    Long-horizon forecasting & MSE & 0.628 & 0.627 & 0.617 \\
    Classification & Accuracy & 0.748 & 0.748 & 0.794 \\
    % Short-horizon forecasting (within-group) & sMAPE & 16.200 & 16.488 & -- \\
    % Short-horizon forecasting (transfer) & sMAPE & 27.210 & 27.368 & 14.525 \\
    Short-horizon forecasting (within-group) & sMAPE & 16.199 & 16.488 & -- \\
    Short-horizon forecasting (transfer) & sMAPE & 27.210 & 27.369 & 14.525 \\
    \bottomrule
  \end{tabular}
\end{table}



\begin{figure}[H]
    \centering
    \includegraphics[width=1\linewidth]{figures/long_horizon_mse.pdf}
    % \caption{Long-horizon forecasting comparison across datasets and forecast horizons.}
    % \caption{Long-horizon forecasting test MSE across datasets and forecast horizons. It compares Raw-$k$NN and PCA-$k$NN, with MOMENT across benchmark settings.}
    \caption{Comparison of long-horizon forecasting using test MSE across datasets and forecast horizons. It compares Raw-$k$NN and PCA-$k$NN with MOMENT across benchmark settings.}
    \label{fig:long_horizon_mse}
\end{figure}



% \begin{figure}[H]
%     \centering
%     \begin{subfigure}[t]{0.48\linewidth}
%         \centering
%         \includegraphics[width=\linewidth]{figures/short_horizon_comparison.pdf}
%         % \caption{Short-horizon forecasting comparison across datasets.}
%         \label{fig:short_horizon_comparison}
%     \end{subfigure}
%     \hfill
%     \begin{subfigure}[t]{0.51\linewidth}
%         \centering
%         \includegraphics[width=\linewidth]{figures/classification_boxplot.pdf}
%         % \caption{Classification accuracy across 91 UCR datasets.}
%         \label{fig:classification_boxplot}
%     \end{subfigure}
%     % \caption{Comparisons on short-horizon forecasting and classification. The left panel shows short-horizon forecasting comparison under within-group and source-to-target transfer settings. The right panel shows classification accuracy across UCR datasets, with MOMENT summary statistics shown as reference lines.}
%     \caption{Comparisons on short-horizon forecasting and classification. The left panel compares forecasting under within-group and source-to-target transfer settings, while the right panel shows classification accuracy across UCR datasets, with MOMENT summary statistics shown as reference lines.}
% \end{figure}

\begin{figure}[H]
    \centering
    \begin{subfigure}[t]{0.48\linewidth}
        \centering
        \includegraphics[width=\linewidth]{figures/short_horizon_comparison.pdf}
        % \caption{Short-horizon forecasting comparison under within-group and source-to-target transfer settings.}
        % \label{fig:short_horizon_comparison}
    \end{subfigure}
    \hfill
    \begin{subfigure}[t]{0.51\linewidth}
        \centering
        \includegraphics[width=\linewidth]{figures/classification_boxplot.pdf}
        % \caption{Classification accuracy across UCR datasets, with MOMENT summary statistics as reference lines.}
        % \label{fig:classification_boxplot}
    \end{subfigure}
    \caption{Comparisons on short-horizon forecasting and classification. The left panel compares forecasting under within-group and source-to-target transfer settings. The right panel shows classification accuracy across UCR datasets, with MOMENT summary statistics shown as reference lines.}
    \label{fig:short_horizon_vs_classification}
\end{figure}

\section{Discussions}\label{sec:discussions}

\paragraph*{Constants versus exponents.}
One main message of the theory is that dependence and local geometry play qualitatively different roles. The local exponent $s$ determines the fundamental scale $(k/n)^{1/s}$, whereas the dependence structure determines constants and concentration quality. This distinction is important both conceptually and empirically: it suggests that weak dependence should affect finite-sample behavior without changing the asymptotic behavior.

\paragraph*{On uniform convergence.} We study nearest-neighbor radii at a fixed query point $x$, both for almost sure convergence and for finite-sample tail and moment bounds. This pointwise analysis matches the local nature of nearest-neighbor methods and yields useful results for downstream learning methods at a given test point \citep{biau2015lectures}. Extending the theory to uniform convergence over the support is interesting, especially for global guarantees on $k$-NN regression, classification, or graph construction \citep{biau2015lectures, jiang2019non}. However, such a result would require additional assumptions and analysis: one would need the lower- and upper-mass conditions to hold uniformly over the support, together with uniform concentration for the dependent counting process indexed by balls, which is typical in empirical process theory under mixing conditions \citep{yu1994rates}. Since our goal is to isolate the effect of dependence on nearest-neighbor radii, we leave the general uniform theory to future work. The pointwise results are also standard in nonparametric theory, where one first identifies the local behavior at a fixed query point and then generalizes to global sup-norm under stronger regularity conditions \citep{gyorfi2002distribution}.

\paragraph*{Implications for learning methods.}
More broadly, we provide a geometric control principle for a wide class of learning methods under dependent sampling. Since the $k$-NN radius determines the effective neighborhood size around a query point, our results imply that, under geometric mixing, dependent observations behave as if the relevant scale were still $(k/n)^{1/s}$ up to constants and a dependence-induced concentration penalty. As a result, many heuristics on the methods from the i.i.d.\ setting continue to apply. In particular, whenever the performance of a method is driven by the size of adaptive neighborhoods---as in $k$-NN regression, local averaging, adaptive bandwidth selection, graph construction, or neighborhood-based manifold methods---our results can be used as an ingredient to transfer the performance to dependent data. From this perspective, we isolate a basic geometric quantity whose behavior under dependence can serve as a building block for future analyses of nonparametric and graph learning algorithms.

\section*{Acknowledgments}
We thank Fang Han for partially motivating this work and for helpful comments.


{%\small
\bibliographystyle{apalike}
\bibliography{AMS}
}

\newpage
\appendix

\section{Proofs}

We begin with a Bernstein-type inequality for \(\alpha\)-mixing sequences.

\begin{lemma}[Theorem 5 of \cite{rio1995functional}; Theorem 2.1 of \cite{liebscher1996strong}]\label{lemma1}
Let \([Z_i]_{i=1}^n\) be a centered sequence of random variables such that
\[
\E[Z_i]=0
\qquad\text{and}\qquad
|Z_i|\le S(n)<\infty,
\qquad i\in \zahl{n}.
\]
Then, for any \(n\ge 1\), any \(m\in \zahl{n}\), and any \(\epsilon>4mS(n)\),
\[
\P\Big(\Big|\sum_{i=1}^n Z_i\Big|>\epsilon\Big)
\le
4\exp \left(
-\frac{\epsilon^2}{64\frac{n}{m}\sigma_{n,m}^2+\frac83\,\epsilon m S(n)}
\right)
+
4\frac{n}{m}\alpha_m,
\]
where
\[
\sigma_{n,m}^2
:=
\sup_{0\le j\le n-1}
\E\Big[\sum_{i=j+1}^{(j+m)\wedge n} Z_i\Big]^2.
\]
\end{lemma}

\subsection{Proof of Theorem~\ref{thm:as}}

\begin{proof}

Fix \(x\in \R^d\). Since \(\mathcal X_{(k)}^n(x)\in {\rm supp}(\mu)\) almost surely and $r_x =\inf\{\|y-x\|: y\in {\rm supp}(\mu)\}$, we have
\[
\|\mathcal X_{(k)}^n(x)-x\| \ge r_x
\qquad\text{a.s. for every }n.
\]
Thus, it remains to prove that
\[
\limsup_{n\to\infty}\|\mathcal X_{(k)}^n(x)-x\| \le r_x
\qquad\text{a.s.}
\]

Fix \(\varepsilon>0\), and define
\[
p_x := \mu(B_{x,r_x+\varepsilon})
= \P\big(\|X_{n,i}-x\|\le r_x+\varepsilon\big).
\]
Since \(B_{x,r_x+\varepsilon}\) intersects \({\rm supp}(\mu)\) by definition of \(r_x\), we have \(p_x>0\).

We observe that the event \(\|\mathcal X_{(k)}^n(x)-x\|>r_x+\varepsilon\) occurs if and only if fewer than \(k\) observations in the $n$-th row fall inside \(B_{x,r_x+\varepsilon}\). Then we have
\[
\P\big(\|\mathcal X_{(k)}^n(x)-x\|>r_x+\varepsilon\big)
=
\P\Big(\sum_{i=1}^n \mathbf 1\{\|X_{n,i}-x\|\le r_x+\varepsilon\}<k\Big).
\]
We set
\[
Z_{n,i}
:=
\mathbf 1\{\|X_{n,i}-x\|\le r_x+\varepsilon\}
-
\P\big(\|X_{n,i}-x\|\le r_x+\varepsilon\big)
=
\mathbf 1\{\|X_{n,i}-x\|\le r_x+\varepsilon\}-p_x.
\]
Then \(\E[Z_{n,i}]=0\) and \(|Z_{n,i}|\le 1\). Since \(Z_{n,i}\) is a measurable function of \(X_{n,i}\), the sequence \([Z_{n,i}]_{i=1}^n\) is again \(\alpha\)-mixing with the same mixing coefficients as \([X_{n,i}]_{i=1}^n\). Moreover, for any \(m\in\{1,\dots,n\}\) and any \(0\le j\le n-1\), the covariance inequality \citep[Corollary of Lemma 2.1]{davydov1968convergence} yields
\begin{align*}
\E\Big[\sum_{i=j+1}^{(j+m)\wedge n} Z_{n,i}\Big]^2
&=
\Var\Big(\sum_{i=j+1}^{(j+m)\wedge n} Z_{n,i}\Big) \\
&\le \sum_{i=j+1}^{(j+m)\wedge n}\Var(Z_{n,i})
+2\sum_{i=j+1}^{(j+m)\wedge n}\sum_{k>i}\Cov(Z_{n,i},Z_{n,k}) \\
&\le
m + 24\sum_{i=j+1}^{(j+m)\wedge n}\sum_{k>i}\alpha_{k-i}
\|Z_{n,i}\|_\infty \|Z_{n,k}\|_\infty \\
&\le
m + 24\sum_{i=j+1}^{(j+m)\wedge n}\sum_{k>i}\alpha_{k-i} \\
&\le
m + 24m\sum_{\ell=1}^\infty \alpha_\ell.
\end{align*}
Since \(\sum_{\ell=1}^\infty \alpha_\ell<\infty\), we set $M = 1+24\sum_{\ell=1}^\infty \alpha_\ell <\infty$, then
\[
\sigma_{n,m}^2
:=
\sup_{0\le j\le n-1}
\E\Big[\sum_{i=j+1}^{(j+m)\wedge n} Z_{n,i}\Big]^2
\le Mm.
\]
Since \(k/n\to 0\) and \(p_x>0\), for all sufficiently large $n$, we have $k-np_x \le -\frac{np_x}{2}$. Then
\begin{align*}
\P\big(\|\mathcal X_{(k)}^n(x)-x\|>r_x+\varepsilon\big)
&=
\P\Big(\sum_{i=1}^n Z_{n,i}<k-np_x\Big) \\
&\le
\P\Big(\Big|\sum_{i=1}^n Z_{n,i}\Big|>np_x-k\Big) \\
&\le
\P\Big(\Big|\sum_{i=1}^n Z_{n,i}\Big|>\frac{np_x}{2}\Big).
\end{align*}

We now apply Lemma~\ref{lemma1} with $m_n = \lfloor \frac{3}{64}\frac{np_x}{\log n}\rfloor$. Then \(m_n\to\infty\) and \(m_n/n \to 0\). Therefore for all sufficiently large $n$, we have $64M \le \frac{4}{3}p_x m_n$ and $
\frac{np_x}{2}>4m_n$, and then the conditions of Lemma~\ref{lemma1} are satisfied with \(S(n)=1\) and \(\epsilon_n=np_x/2\). Using \(\sigma_{n,m_n}^2\le M m_n\), we obtain
\begin{align*}
\P\Big(\Big|\sum_{i=1}^n Z_{n,i}\Big|>\frac{np_x}{2}\Big)
&\le
4\exp \left[
-\frac{(np_x/2)^2}{64(n/m_n)\sigma_{n,m_n}^2 + \frac83 (np_x/2)m_n}
\right]
+4\frac{n}{m_n}\alpha_{m_n} \\
&\le
4\exp \left[
-\frac{n^2p_x^2/4}{64Mn + \frac43 np_x m_n}
\right]
+4\frac{n}{m_n}\alpha_{m_n} \\
&\le
4\exp \left[
-\frac{n^2p_x^2/4}{\frac83 np_x m_n}
\right]
+4\frac{n}{m_n}\alpha_{m_n} \\
&=
4\exp \left(-\frac{3}{32}\frac{np_x}{m_n}\right)
+4\frac{n}{m_n}\alpha_{m_n}.
\end{align*}
By the choice of \(m_n\), the exponential term is \(O(n^{-2})\).

We now assume the mixing coefficients satisfy $\alpha_m \lesssim m^{-(1+\gamma)}$ for some $\gamma>1$. Then
\[
\frac{n}{m_n}\alpha_{m_n} \lesssim \frac{n}{m_n^{2+\gamma}} \lesssim \frac{(\log n)^{2+\gamma}}{n^{1+\gamma}},
\]
which is summable in $n$. Then
\[
\sum_{n=1}^\infty
\P\big(\|\mathcal X_{(k)}^n(x)-x\|>r_x+\varepsilon\big)
<\infty.
\]
By the Borel--Cantelli lemma,
\[
\P\big(\|\mathcal X_{(k)}^n(x)-x\|>r_x+\varepsilon\ \text{i.o.}\big)=0.
\]
That is, for every fixed \(\varepsilon>0\),
\[
\|\mathcal X_{(k)}^n(x)-x\|\le r_x+\varepsilon
\qquad\text{eventually a.s.}
\]

Let \(\varepsilon=1/j\), \(j\in\mathbb N\), and intersect the resulting probability-one events. Then we obtain
\[
\limsup_{n\to\infty}\|\mathcal X_{(k)}^n(x)-x\|\le r_x
\qquad\text{a.s.}
\]
Together with the lower bound \(\|\mathcal X_{(k)}^n(x)-x\|\ge r_x\) a.s. for every $n$, we obtain
\[
\|\mathcal X_{(k)}^n(x)-x\| \to r_x
\qquad\text{almost surely}.
\]
This completes the proof.

\end{proof}

\subsection{Proof of Theorem~\ref{thm:tail}}

\begin{proof}

Fix \(r>0\), and define
\[
Y_i(r):=\mathbf 1\{\|X_{n,i}-x\|\le r\},
\qquad
p_r:=\E[Y_i(r)]=\mu(B(x,r)).
\]
Let
\[
N_n(x,r):=\sum_{i=1}^n Y_i(r).
\]
By the definition of the \(k\)-nearest-neighbor radius,
\[
R_{n,k}(x)>r
\iff N_n(x,r)<k.
\]
We now use a blocking argument. Partition the row into alternating kept and discarded blocks of length \(b_n\). Set $M_n=\lfloor \frac{n}{2b_n}\rfloor$, and for \(u=1,\dots,M_n\), define the kept block
\[
I_u:=\{(2u-2)b_n+1,\dots,(2u-1)b_n\}.
\]
Let
\[
W_u(r):=\sum_{i\in I_u}Y_i(r),
\qquad
S_n^{\rm keep}(r):=\sum_{u=1}^{M_n}W_u(r).
\]
Since \(S_n^{\rm keep}(r)\le N_n(x,r)\), we obtain
\[
\mathbb P(R_{n,k}(x)>r)
=
\mathbb P(N_n(x,r)<k)
\le
\mathbb P(S_n^{\rm keep}(r)<k).
\]

Consider the radius
\[
r_j=2^j r_*(n,k)
=
2^j\Big(\frac{8k}{c_- n}\Big)^{1/s},
\qquad j\ge 0,
\]
and assume \(r_j\le r_0\). By Assumption~\ref{ass:lower},
\[
p_{r_j}
=
\mu(B(x,r_j))
\ge c_- r_j^s
=
8\cdot2^{js}\frac{k}{n}.
\]
Since $b_n/n \to 0$ and $M_n b_n \ge n/4$ for all sufficiently large $n$,
\[
\mu_j=\E[S_n^{\rm keep}(r_j)]
=
M_n b_n p_{r_j}
\ge n\,p_{r_j}/4
\ge
2^{js+1}k.
\]
Therefore, on the event \(\{S_n^{\rm keep}(r_j)<k\}\), since $\mu_j\ge 2k$,
\[
\mu_j-S_n^{\rm keep}(r_j)>\mu_j-k\ge \mu_j/2.
\]
Then
\[
\{S_n^{\rm keep}(r_j)<k\}
\subseteq
\left\{
\left|S_n^{\rm keep}(r_j)-\mu_j\right|>\frac{\mu_j}{2}
\right\}.
\]
Therefore, if we define
\[
Z_u:=W_u(r_j)-\E[W_u(r_j)],
\qquad u=1,\dots,M_n,
\]
we have
\[
\mathbb P(R_{n,k}(x)>r_j)
\le
\mathbb P\Big(\Big|\sum_{u=1}^{M_n}Z_u\Big|>\frac{\mu_j}{2}\Big).
\]

Next we apply Lemma~\ref{lemma1} for the block sequence \([Z_u]_{u=1}^{M_n}\). Since each \(Z_u\) is measurable with respect to the coordinates in the kept block \(I_u\), and successive kept blocks are separated by a discarded block of length \(b_n\), the sequence \([Z_u]_{u=1}^{M_n}\) is \(\alpha\)-mixing with coefficients
\[
\bar\alpha_\ell
:=
\sup_{u\ge 1}\sup_{A\in\sigma(Z_1,\dots,Z_u),B\in\sigma(Z_{u+\ell},Z_{u+\ell+1},\dots)}
|\mathbb P(A\cap B)-\mathbb P(A)\mathbb P(B)|
\le
\alpha_{\ell b_n},
\qquad \ell\ge 1.
\]
In particular, $\bar\alpha_1\le \alpha_{b_n}\le e^{-\gamma b_n}\le n^{-8}$, as long as \(b_n\) is chosen with a sufficiently large constant times \(\log n\). Also, $0\le W_u(r_j)\le b_n$, so $|Z_u| = |W_u(r_j)-\E W_u(r_j)| \le b_n$. Thus the boundedness parameter in Lemma~\ref{lemma1} is $S(M_n)=b_n$.

To bound the variance term \(\sigma^2_{M_n,1}\), since \(W_u(r_j)\) is a sum of \(b_n\) Bernoulli indicators, we have $W_u(r_j)^2\le b_n W_u(r_j)$, and therefore
\[
\Var(Z_u)
\le \E[W_u(r_j)^2]
\le b_n\,\E[W_u(r_j)]
=
b_n^2 p_{r_j}.
\]
Then for the block sequence with block parameter \(m=1\),
\[
\sigma^2_{M_n,1}
=
\sup_{0\le v\le M_n-1}\E[Z_{v+1}^2]
\le
b_n^2 p_{r_j}.
\]

We now apply Lemma~\ref{lemma1} to \([Z_u]_{u=1}^{M_n}\) with \(m=1\) and threshold $\varepsilon=\mu_j/2=M_n b_n p_{r_j}/2$. Since $\varepsilon\ge k\ge K_0\log n$, and \(b_n\asymp \log n\), we may choose \(K_0\) sufficiently large so that $\varepsilon>4b_n$ for all sufficiently large $n$. Then by Lemma~\ref{lemma1},
\begin{align*}
\mathbb P\Big(\Big|\sum_{u=1}^{M_n}Z_u\Big|>\frac{\mu_j}{2}\Big)
&\le
4\exp \left(
-\frac{(\mu_j/2)^2}
{64 M_n \sigma_{M_n,1}^2 + \frac83(\mu_j/2)b_n}
\right)
+4M_n\bar\alpha_1 \\
&\le
4\exp \left(
-\frac{(\mu_j/2)^2}
{64 M_n b_n^2 p_{r_j} + \frac43 \mu_j b_n}
\right)
+4M_n\bar\alpha_1.
\end{align*}
Using \(\mu_j=M_n b_n p_{r_j}\), we obtain $\frac43 \mu_j b_n
=
\frac43 M_n b_n^2 p_{r_j}$, so the denominator is bounded by
\[
\left(64+\frac43\right)M_n b_n^2 p_{r_j}
\le 66\,M_n b_n^2 p_{r_j},
\]
and then
\begin{align*}
\frac{(\mu_j/2)^2}
{64 M_n b_n^2 p_{r_j} + \frac43 \mu_j b_n}
&\ge
c\frac{\mu_j^2}{M_n b_n^2 p_{r_j}} =
c\frac{(M_n b_n p_{r_j})^2}{M_n b_n^2 p_{r_j}} =
c M_n p_{r_j}.
\end{align*}
Since \(M_n\ge n/(4b_n)\) for all large $n$,
\[
M_n p_{r_j}
\ge
\frac{n}{4b_n}\cdot 8\cdot2^{js}\frac{k}{n}
=
2\cdot2^{js}\frac{k}{b_n}.
\]
Therefore, for some universal constant \(c_0>0\),
\[
\mathbb P(R_{n,k}(x)>r_j)
\le
4\exp \left(-c_0 2^{js}\frac{k}{b_n}\right)
+
4M_n\bar\alpha_1.
\]

Finally, using \(b_n\asymp \log n\), \(\bar\alpha_1\le n^{-8}\), and \(M_n\le n/(2b_n)\), we obtain $4M_n\bar\alpha_1 \le C n^{-7}$, and then
\[
\mathbb P(R_{n,k}(x)>r_j)
\le
4\exp \left(-c_0\,2^{js}\frac{k}{\log n}\right)
+
C n^{-7},
\]
after adjusting constants. This completes the proof.
\end{proof}


\subsection{Proof of Theorem~\ref{thm:moment}}

\begin{proof}

Let $$r_* = \Big(\frac{8k}{c_- n}\Big)^{1/s}$$ and choose \(J_0\) maximal such that $2^{J_0}r_* \le r_0$. Since the support of \(\mu\) has diameter at most \(D\), we have $R_{n,k}(x)\le D$ almost surely.

We decompose the moment over
\[
\{R_{n,k}(x)\le r_*\},
\quad
\{2^j r_* < R_{n,k}(x)\le 2^{j+1}r_*\},\ 0\le j\le J_0-1,
\quad
\{R_{n,k}(x)>2^{J_0}r_*\},
\]
and we obtain
\begin{align*}
\mathbb E[R_{n,k}(x)^p]
&\le r_*^p
+ \sum_{j=0}^{J_0-1}(2^{j+1}r_*)^p\,\mathbb P(R_{n,k}(x)>2^j r_*) + D^p\,\mathbb P(R_{n,k}(x)>2^{J_0}r_*).
\end{align*}

We now apply Theorem~\ref{thm:tail}. For every \(0\le j\le J_0\),
\[
\mathbb P(R_{n,k}(x)>2^j r_*)
\le
4\exp \Big(-c_0\,2^{js}\frac{k}{\log n}\Big)+Cn^{-7},
\]
and then
\begin{align*}
\mathbb E[R_{n,k}(x)^p]
&\le r_*^p
+ C r_*^p \sum_{j=0}^{J_0-1}2^{jp}
\exp \Big(-c_0\,2^{js}\frac{k}{\log n}\Big) + C r_*^p n^{-7}\sum_{j=0}^{J_0-1}2^{jp}
+ C D^p n^{-7}.
\end{align*}

We now bound the three terms on the right-hand side. For the first term, since \(k\ge K_0\log n\), the quantity \(k/\log n\) is bounded below by a positive constant. Therefore
\[
\sum_{j=0}^{J_0-1}2^{jp}
\exp \Big(-c_0\,2^{js}\frac{k}{\log n}\Big)
\le
\sum_{j\ge 0}2^{jp}e^{-c_1 2^{js}}
<\infty
\]
for some constant \(c_1>0\). For the second term, since \(2^{J_0}r_*\le r_0<2^{J_0+1}r_*\), we have $2^{J_0}\asymp r_*^{-1}$, where the constants depend only on \(r_0\). Then
\[
r_*^p\sum_{j=0}^{J_0-1}2^{jp}
\le
C r_*^p 2^{J_0 p}
\le C,
\]
and then
\[
C r_*^p n^{-7}\sum_{j=0}^{J_0-1}2^{jp}
\le C n^{-7},
\]
which is negligible. The third term \(C D^p n^{-7}\) is negligible as well. Combining the above estimates yields
\[
\mathbb E[R_{n,k}(x)^p]
\le C r_*^p + C n^{-7}
\le C r_*^p.
\]
Since
\[
r_*^p
=
\Big(\frac{8k}{c_- n}\Big)^{p/s}
\asymp
\Big(\frac{k}{n}\Big)^{p/s},
\]
we obtain
\[
\mathbb E[R_{n,k}(x)^p]
\le
C\Big(\frac{k}{n}\Big)^{p/s}.
\]
This completes the proof.
\end{proof}


\subsection{Proof of Theorem~\ref{thm:lower}}

\begin{proof}
Let $$r^- = \Big(\frac{k}{2c_+ n}\Big)^{1/s}.$$ By Assumption~\ref{ass:upper},
\[
\mu(B(x,r^-)) \le c_+(r^-)^s = \frac{k}{2n}.
\]
Therefore,
\[
\E[N_n(x,r^-)]
=
n\,\mu(B(x,r^-))
\le k/2,
\]
where,
\[
N_n(x,r)=\sum_{i=1}^n \mathbf 1\{\|X_{n,i}-x\|\le r\}.
\]

By Markov's inequality,
\[
\P\big(N_n(x,r^-)\ge k\big)
\le
\frac{\E[N_n(x,r^-)]}{k}
\le \frac12,
\]
and
\[
R_{n,k}(x)>r^-
\iff
N_n(x,r^-)<k,
\]
we have
\[
\P\big(R_{n,k}(x)>r^-\big)\ge \frac12.
\]

On the event \(\{R_{n,k}(x)>r^-\}\), we have \(R_{n,k}(x)^p\ge (r^-)^p\), and then
\[
\E[R_{n,k}(x)^p]
\ge
(r^-)^p\,\P\big(R_{n,k}(x)>r^-\big)
\ge
\frac12\,(r^-)^p.
\]
By the definition of \(r^-\), we obtain
\[
\E[R_{n,k}(x)^p]
\ge
\frac12\Big(\frac{k}{2c_+ n}\Big)^{p/s}.
\]
The completes the proof.
\end{proof}

\subsection{Proof of Theorem~\ref{thm:exact}}

\begin{proof}
The conclusion follows directly by combining the upper and lower bounds established in Theorems~\ref{thm:moment} and \ref{thm:lower}.
\end{proof}


\section{Experiments}
We provide additional details on the implementation of the experiments.

\subsection{Experiment 1 (Section~\ref{exp1:Moment convergence rate}) implementation details}
\label{app: Experiment 1 implementation details}

\paragraph*{Setup.}
Experiment 1 uses interior query aggregation. For each MC replication, we sample 
\[
X_{\mathrm{eval}}\sim \mathrm{Unif}([0.01,0.99]^d),\qquad |X_{\mathrm{eval}}|=1000,
\]
The number of NNs is chosen as $k=\lfloor n^\beta\rceil$ with $\beta$ on a 20-point grid in $[0.1,0.9]$ for slope fitting. We also impose a cap for $\frac{k}{n}\le 0.01$. 

\paragraph*{Data generating processes.}
The i.i.d. uniform data is used as the baseline in every panel. The dependent data is generated by the following processes, with three strength levels, as shown in Table~\ref{tab:exp1-dgp-settings}.

\begin{table}
  \caption{Dependence-strength settings for the three dependent DGPs in experiment 1. Weak settings are near i.i.d., while medium and strong settings increase dependence contrast.}
  \label{tab:exp1-dgp-settings}
  \centering
  \small
  \begin{tabular}{@{}lllll@{}}
    \toprule
    Model & Tuning parameter & Weak & Medium & Strong \\
    \midrule
    Linear Gaussian state-space (factor AR(1)) & AR coefficient $\rho$ & $0.02$ & $0.60$ & $0.98$ \\
    Hidden Markov model (2-state) & Stay probability $p_{\mathrm{stay}}$ & $0.70$ & $0.99$ & $0.9995$ \\
    Gaussian process (FFT, RBF) & Lengthscale $\ell$ & $1$ & $50$ & $200$ \\
    \bottomrule
  \end{tabular}
\end{table}


\begin{itemize}
\item \textbf{Linear Gaussian state-space (LSS).}
A shared latent AR(1) factor is generated as
\[
f_t=\rho f_{t-1}+\sqrt{1-\rho^2}\,\eta_t,\qquad \eta_t\sim\mathcal N(0,1).
\]
Then each coordinate is
\[
z_{t,j}=a f_t+\sqrt{1-a^2}\,\varepsilon_{t,j},\qquad \varepsilon_{t,j}\sim\mathcal N(0,1),
\]
with dimension-damped shared loading
\[
a=\frac{s\sqrt{|\rho|}}{d^\gamma},\qquad s=0.9,\ \gamma=0.3.
\]
(Implementation clips $a\in[0,0.999]$ and uses burn-in 400.)

\item \textbf{Hidden Markov model (HMM).}
Each coordinate $j$ has its own two-state chain $S_{t,j}\in\{0,1\}$ with
\[
\Pr(S_{t,j}=S_{t-1,j})=p_{\mathrm{stay}}.
\]
Conditional emissions are
\[
Z_{t,j}\mid S_{t,j}\sim \mathcal N \big((2S_{t,j}-1)\mu_{\mathrm{eff}},\,\sigma^2\big),
\]
where $\mu=1.75$, $\sigma=0.70$, and
\[
\mu_{\mathrm{eff}}=\mu\,|2p_{\mathrm{stay}}-1|^6.
\]
(This keeps weak settings close to i.i.d.\ after marginal transformation while separating medium/strong regimes; burn-in 400.)

\item \textbf{Gaussian process (FFT, RBF).}
For each coordinate, a stationary Gaussian sequence is generated with lag covariance
\[
k(h)=\exp \left(-\frac{h^2}{2\ell^2}\right),\qquad \ell\in\{1,50,200\},
\]
using circulant embedding and FFT spectral filtering ($O(n\log n)$), then truncating to length $n$.

\item \textbf{Transformation.}
All $d$ coordinates are transformed to have approximately $\mathrm{Unif}[0,1]$ marginal distributions while retaining dependence:
\[
U_{t,j}=\Phi(Z_{t,j}) \quad \text{(LSS/GP)},
\]
and, for HMM, via the stationary two-component mixture CDF
\[
F_{\mathrm{mix}}(z)=\tfrac12\Phi \left(\frac{z+\mu_{\mathrm{eff}}}{\sigma}\right)
+\tfrac12\Phi \left(\frac{z-\mu_{\mathrm{eff}}}{\sigma}\right).
\]
\end{itemize}

\subsection{Experiment 2 (Section~\ref{exp2: Entropy estimation}) implementation details}
\label{app: Experiment 2 implementation details}

\paragraph*{Data-generating process.}
For each intrinsic dimension $s$, we generate independent stationary Gaussian AR(1) latent processes across coordinates:
\[
Z_t=\rho Z_{t-1}+\sqrt{1-\rho^2}\,\varepsilon_t,
\qquad \varepsilon_t\sim \mathcal N(0,I_s).
\]
The temporal dependence levels are $\rho=0$ for the i.i.d.\ benchmark, $\rho=0.3$ for weak AR(1) dependence, and $\rho=0.6$ for strong AR(1) dependence. The latent coordinates are independent across dimensions, while each coordinate follows the same AR(1) dependence over time. The latent vector is embedded into a $5$-dimensional signal representation and then mapped into $20$ observed coordinates by a fixed orthonormal $20\times 5$ linear map, which induces a low-rank cross-coordinate dependence structure in the ambient observations. The PCA estimator fits a $5$-component PCA representation to the observed data and then estimates entropy using the top $s$ PCA coordinates, which, in this linear Gaussian setting, recover an orthogonal low-dimensional representation up to rotation.

\paragraph*{Estimators and simulation grid.}
All reported curves use the Kozachenko--Leonenko $k$-NN entropy estimator with Euclidean distance. The oracle intrinsic estimate is computed on the true latent $s$-dimensional sample, and the PCA estimate is computed on the top $s$ PCA coordinates. The current implementation uses $1000$ Monte Carlo replications with global random seed 20260504, burn-in $2000$, and neighbor schedule $k=\max\{2,\lceil n^{0.1}\rceil\}$.

\subsection{Experiment 3 (Section~\ref{exp3: Real-data_evaluation}) implementation details}
\label{Appendix: real_data_implementation details}

We follow the real-data benchmark setup of \cite{goswami2024moment} for long-horizon forecasting, short-horizon forecasting, and classification, using the same task definitions, evaluation metrics, and preprocessing procedures. In addition to the source-to-target short-horizon setting considered in \cite{goswami2024moment}, we also report a within-group short-horizon evaluation. The setup for each task is as follows.

\begin{itemize}
    \item \textbf{Long-horizon forecasting.} We use 9 datasets from the Informer long-horizon forecasting benchmark: ETTm1, ETTm2, ETTh1, ETTh2, Electricity, Traffic, Weather, Exchange, and ILI. The evaluation metrics are MSE and mean absolute error (MAE). We use a look-back window length of $512$ and a horizontal 60/10/30 train/validation/test split, with forecast horizons $\{24,36,48,60\}$ for ILI and $\{96,192,336,720\}$ otherwise. Hyperparameters are selected by 5-fold time-series cross-validation when the training set contains at least 20 windows, and by validation-holdout tuning otherwise.

    \item \textbf{Short-horizon forecasting.} We use 6 benchmark groups from the Monash time series forecasting archive: M4-Yearly, M4-Quarterly, M4-Monthly, M3-Yearly, M3-Quarterly, and M3-Monthly. The evaluation metric is sMAPE, and the input length is $512$. In the source-to-target transfer setting, hyperparameters are selected on a source group and evaluated directly on a different target group with the same frequency. We also report a within-group short-horizon evaluation, which is not included in the original MOMENT benchmark. In this setting, each dataset group is randomly split into a 70\% source portion and a 30\% target portion; hyperparameters are selected by 5-fold cross-validation on the source portion, and evaluation is performed on the target portion from the same group.

    \item \textbf{Classification.} We use 91 datasets from the UCR classification archive. The evaluation metric is accuracy. Inputs are represented as length-$512$ time series, and we use the official train/test splits provided by the UCR archive. Hyperparameters are selected using only the official training split, via stratified cross-validation when feasible, and final performance is evaluated on the official test split.
\end{itemize}


% \begin{table}[H]
% \centering
% \caption{Summary of the real-data benchmark tasks, datasets, evaluation metrics, and experimental settings. We follow the benchmark protocol of \citep{goswami2024moment} while evaluating Raw-$k$NN and PCA-$k$NN.}
% \label{tab:real_data_benchmark_setup}
% \small
% \begin{tabular}{C{2.8cm}|C{3.8cm}|C{2.2cm}|L{5.8cm}}
% \hline
% \textbf{Task} & \textbf{Datasets} & \textbf{Metrics} & \textbf{Experimental setting} \\
% \hline
% Long-horizon forecasting
% & 9 datasets from the Informer long-horizon forecasting benchmark: ETTm1, ETTm2, ETTh1, ETTh2, Electricity, Traffic, Weather, Exchange, ILI
% & MSE, MAE
% & Look-back window length $512$; horizontal split; horizons $\{24,60\}$ for ILI and $\{96,720\}$ otherwise; predefined splits when available, otherwise 60/10/30; 5-fold time-series CV or validation-holdout tuning. \\
% \hline
% Short-horizon forecasting
% & 6 benchmark groups from the Monash time series forecasting archive: M4-Yearly, M4-Quarterly, M4-Monthly, M3-Yearly, M3-Quarterly, M3-Monthly
% & sMAPE
% & Input length $512$; source-to-target transfer across matching frequency groups; source-only CV and direct target evaluation. \\
% \hline
% Classification
% & 91 datasets from the UCR classification archive
% & Accuracy
% & Inputs represented as length-$512$ time series; official train/test splits; model selection on the training split only. \\
% % \hline
% % Imputation
% % & 6 datasets from the Informer long-horizon forecasting benchmark: ETTm1, ETTm2, ETTh1, ETTh2, Electricity, Weather
% % & MSE, MAE
% % & Input length $512$; same forecasting splits; patch length $8$; masking ratios $\{12.5\%,25\%,37.5\%,50\%\}$; masked validation tuning. \\
% \hline
% \end{tabular}
% \end{table}


\begin{figure}[H]
    \centering
    \includegraphics[width=1\linewidth]{figures/long_horizon_mae.pdf}
    % \caption{Long-horizon forecasting comparison across datasets and forecast horizons in terms of MAE. }
    % \caption{Long-horizon forecasting test MAE across datasets and forecast horizons. It compares nearest-neighbor baselines, Raw-$k$NN and PCA-$k$NN, with MOMENT in several settings.}
    % \caption{Long-horizon forecasting test MAE across datasets and forecast horizons. It compares Raw-$k$NN and PCA-$k$NN with MOMENT across benchmark settings.}
    \caption{Comparison of long-horizon forecasting using test MAE across datasets and forecast horizons. It compares Raw-$k$NN and PCA-$k$NN with MOMENT across benchmark settings.}
    \label{fig:long_horizon_mae}
\end{figure}

% \begin{figure}[H]
%     \centering
%     \includegraphics[width=0.72\linewidth]{figures/imputation_mae.pdf}
%     \caption{Imputation comparison across datasets in terms of MAE. Lower values are better.}
%     \label{fig:imputation_mae}
% \end{figure}



% \begin{table}[H]
% \centering
% \caption{Long-horizon forecasting results. We report test MSE and MAE for Raw-$k$NN, PCA-$k$NN, and MOMENT, with MOMENT results taken from Table 2 of \citep{goswami2024moment}.}
% \label{tab:long_horizon_results}
% \small
% {\footnotesize
% \begin{tabular}{l|l|c|c|c|c|c|c}
% \hline
% \multicolumn{2}{c|}{\textbf{Methods}} & \multicolumn{2}{c|}{\textbf{Raw-$k$NN}} & \multicolumn{2}{c|}{\textbf{PCA-$k$NN}} & \multicolumn{2}{c}{\textbf{MOMENT$_{\mathrm{LP}}$}} \\
% \hline
% \textbf{Dataset} & \textbf{Horizon} & \textbf{MSE} & \textbf{MAE} & \textbf{MSE} & \textbf{MAE} & \textbf{MSE} & \textbf{MAE} \\
% \hline
% \multirow{4}{*}{Weather}
% & 96  & 0.023 & 0.036 & 0.121 & 0.043 & 0.154 & 0.209 \\
% & 192 & 0.002 & 0.034 & 0.003 & 0.037 & 0.197 & 0.248 \\
% & 336 & 0.055 & 0.052 & 0.160 & 0.047 & 0.246 & 0.285 \\
% & 720 & 0.095 & 0.068 & 0.103 & 0.064 & 0.315 & 0.336 \\
% \hline
% \multirow{4}{*}{ETTh1}
% & 96  & 0.287 & 0.448 & 0.281 & 0.443 & 0.387 & 0.410 \\
% & 192 & 0.237 & 0.407 & 0.227 & 0.397 & 0.410 & 0.426 \\
% & 336 & 0.239 & 0.408 & 0.229 & 0.401 & 0.422 & 0.437 \\
% & 720 & 0.305 & 0.461 & 0.294 & 0.453 & 0.454 & 0.472 \\
% \hline
% \multirow{4}{*}{ETTh2}
% & 96  & 0.289 & 0.425 & 0.286 & 0.423 & 0.288 & 0.345 \\
% & 192 & 0.306 & 0.436 & 0.309 & 0.444 & 0.349 & 0.386 \\
% & 336 & 0.332 & 0.456 & 0.318 & 0.453 & 0.369 & 0.408 \\
% & 720 & 0.433 & 0.525 & 0.410 & 0.515 & 0.403 & 0.439 \\
% \hline
% \multirow{4}{*}{ETTm1}
% & 96  & 0.314 & 0.465 & 0.297 & 0.455 & 0.293 & 0.349 \\
% & 192 & 0.382 & 0.520 & 0.356 & 0.502 & 0.326 & 0.368 \\
% & 336 & 0.325 & 0.481 & 0.328 & 0.484 & 0.352 & 0.384 \\
% & 720 & 0.304 & 0.461 & 0.247 & 0.416 & 0.405 & 0.416 \\
% \hline
% \multirow{4}{*}{ETTm2}
% & 96  & 0.212 & 0.356 & 0.206 & 0.352 & 0.170 & 0.260 \\
% & 192 & 0.245 & 0.387 & 0.215 & 0.361 & 0.227 & 0.297 \\
% & 336 & 0.252 & 0.393 & 0.246 & 0.389 & 0.275 & 0.328 \\
% & 720 & 0.279 & 0.415 & 0.275 & 0.414 & 0.363 & 0.387 \\
% \hline
% \multirow{4}{*}{ILI}
% & 24 & 3.038 & 1.306 & 3.065 & 1.313 & 2.728 & 1.114 \\
% & 36 & 2.723 & 1.245 & 2.728 & 1.244 & 2.669 & 1.092 \\
% & 48 & 2.915 & 1.263 & 2.848 & 1.279 & 2.728 & 1.098 \\
% & 60 & 2.872 & 1.265 & 2.931 & 1.316 & 2.883 & 1.126 \\
% \hline
% \multirow{4}{*}{ECL}
% & 96  & 0.476 & 0.508 & 0.470 & 0.508 & 0.136 & 0.233 \\
% & 192 & 0.493 & 0.518 & 0.493 & 0.522 & 0.152 & 0.247 \\
% & 336 & 0.509 & 0.528 & 0.502 & 0.529 & 0.167 & 0.264 \\
% & 720 & 0.545 & 0.555 & 0.520 & 0.544 & 0.205 & 0.295 \\
% \hline
% \multirow{4}{*}{Traffic}
% & 96  & 0.396 & 0.447 & 0.392 & 0.449 & 0.391 & 0.282 \\
% & 192 & 0.405 & 0.453 & 0.400 & 0.454 & 0.404 & 0.287 \\
% & 336 & 0.408 & 0.457 & 0.405 & 0.458 & 0.414 & 0.292 \\
% & 720 & 0.409 & 0.462 & 0.410 & 0.463 & 0.450 & 0.310 \\
% \hline
% \end{tabular}
% }
% \end{table}

\begin{table}
  \caption{Long-horizon forecasting results. We report test MSE and MAE for Raw-$k$NN, PCA-$k$NN, and MOMENT; MOMENT results are taken from Table 2 of \cite{goswami2024moment}.}
  \label{tab:long_horizon_results}
  \centering
  \footnotesize
  \begin{tabular}{@{}lccccccc@{}}
    \toprule
    \multicolumn{2}{c}{Setting} &
    \multicolumn{2}{c}{Raw-$k$NN} &
    \multicolumn{2}{c}{PCA-$k$NN} &
    \multicolumn{2}{c}{MOMENT$_{\mathrm{LP}}$} \\
    \cmidrule(r){1-2}
    \cmidrule(r){3-4}
    \cmidrule(r){5-6}
    \cmidrule(l){7-8}
    Dataset & Horizon & MSE & MAE & MSE & MAE & MSE & MAE \\
    \midrule
    \multirow{4}{*}{Weather}
      & 96  & 0.023 & 0.036 & 0.121 & 0.043 & 0.154 & 0.209 \\
      & 192 & 0.002 & 0.034 & 0.003 & 0.037 & 0.197 & 0.248 \\
      & 336 & 0.055 & 0.052 & 0.160 & 0.047 & 0.246 & 0.285 \\
      & 720 & 0.095 & 0.068 & 0.103 & 0.064 & 0.315 & 0.336 \\
    \midrule
    \multirow{4}{*}{ETTh1}
      & 96  & 0.287 & 0.448 & 0.281 & 0.443 & 0.387 & 0.410 \\
      & 192 & 0.237 & 0.407 & 0.227 & 0.397 & 0.410 & 0.426 \\
      & 336 & 0.239 & 0.408 & 0.229 & 0.401 & 0.422 & 0.437 \\
      & 720 & 0.305 & 0.461 & 0.294 & 0.453 & 0.454 & 0.472 \\
    \midrule
    \multirow{4}{*}{ETTh2}
      & 96  & 0.289 & 0.425 & 0.286 & 0.423 & 0.288 & 0.345 \\
      & 192 & 0.306 & 0.436 & 0.309 & 0.444 & 0.349 & 0.386 \\
      & 336 & 0.332 & 0.456 & 0.318 & 0.453 & 0.369 & 0.408 \\
      & 720 & 0.433 & 0.525 & 0.410 & 0.515 & 0.403 & 0.439 \\
    \midrule
    \multirow{4}{*}{ETTm1}
      & 96  & 0.314 & 0.465 & 0.297 & 0.455 & 0.293 & 0.349 \\
      & 192 & 0.382 & 0.520 & 0.356 & 0.502 & 0.326 & 0.368 \\
      & 336 & 0.325 & 0.481 & 0.328 & 0.484 & 0.352 & 0.384 \\
      & 720 & 0.304 & 0.461 & 0.247 & 0.416 & 0.405 & 0.416 \\
    \midrule
    \multirow{4}{*}{ETTm2}
      & 96  & 0.212 & 0.356 & 0.206 & 0.352 & 0.170 & 0.260 \\
      & 192 & 0.245 & 0.387 & 0.215 & 0.361 & 0.227 & 0.297 \\
      & 336 & 0.252 & 0.393 & 0.246 & 0.389 & 0.275 & 0.328 \\
      & 720 & 0.279 & 0.415 & 0.275 & 0.414 & 0.363 & 0.387 \\
    \midrule
    \multirow{4}{*}{ILI}
      & 24 & 3.038 & 1.306 & 3.065 & 1.313 & 2.728 & 1.114 \\
      & 36 & 2.723 & 1.245 & 2.728 & 1.244 & 2.669 & 1.092 \\
      & 48 & 2.915 & 1.263 & 2.848 & 1.279 & 2.728 & 1.098 \\
      & 60 & 2.872 & 1.265 & 2.931 & 1.316 & 2.883 & 1.126 \\
    \midrule
    \multirow{4}{*}{ECL}
      & 96  & 0.476 & 0.508 & 0.470 & 0.508 & 0.136 & 0.233 \\
      & 192 & 0.493 & 0.518 & 0.493 & 0.522 & 0.152 & 0.247 \\
      & 336 & 0.509 & 0.528 & 0.502 & 0.529 & 0.167 & 0.264 \\
      & 720 & 0.545 & 0.555 & 0.520 & 0.544 & 0.205 & 0.295 \\
    \midrule
    \multirow{4}{*}{Traffic}
      & 96  & 0.396 & 0.447 & 0.392 & 0.449 & 0.391 & 0.282 \\
      & 192 & 0.405 & 0.453 & 0.400 & 0.454 & 0.404 & 0.287 \\
      & 336 & 0.408 & 0.457 & 0.405 & 0.458 & 0.414 & 0.292 \\
      & 720 & 0.409 & 0.462 & 0.410 & 0.463 & 0.450 & 0.310 \\
    \bottomrule
  \end{tabular}
\end{table}




% \begin{table}[H]
% \centering
% \caption{Classification accuracy across 91 UCR datasets. We report summary statistics for Raw-$k$NN, PCA-$k$NN, and MOMENT. MOMENT results are taken from \citep{goswami2024moment} in Table 4.}
% \label{tab:classification_results}
% \small
% \begin{tabular}{l|c|c|c}
% \hline
%  & \textbf{Raw-$k$NN} & \textbf{PCA-$k$NN} & \textbf{MOMENT$_0$} \\
% \hline
% Mean    & 0.748 & 0.748 & 0.794 \\
% Median  & 0.753 & 0.760 & 0.815 \\
% Std.    & 0.153 & 0.151 & 0.147 \\
% \hline
% \end{tabular}
% \end{table}

\begin{table}
  \caption{Classification accuracy across 91 UCR datasets. We report summary statistics for Raw-$k$NN, PCA-$k$NN, and MOMENT; MOMENT results are taken from Table 4 of \cite{goswami2024moment}.}
  \label{tab:classification_results}
  \centering
  \small
  \begin{tabular}{@{}lccc@{}}
    \toprule
    Statistic & Raw-$k$NN & PCA-$k$NN & MOMENT$_0$ \\
    \midrule
    Mean   & 0.748 & 0.748 & 0.794 \\
    Median & 0.753 & 0.760 & 0.815 \\
    Std.   & 0.153 & 0.151 & 0.147 \\
    \bottomrule
  \end{tabular}
\end{table}


% \begin{table}[H]
% \centering
% \caption{Short-horizon forecasting results measured using sMAPE. We report our Raw-$k$NN and PCA-$k$NN transfer results together with MOMENT results from \citep{goswami2024moment} in Table 3 under the same source-to-target evaluation setting.}
% \label{tab:short_horizon_results_transfer}
% \small
% \begin{tabular}{l|l|c|c|c|c}
% \hline
% \multicolumn{2}{c|}{\multirow{2}{*}{\textbf{Dataset}}} & \multirow{2}{*}{\textbf{Raw-$k$NN}} & \multirow{2}{*}{\textbf{PCA-$k$NN}} & \multicolumn{2}{c}{\textbf{MOMENT$_{\mathrm{LP}}$}} \\
% \cline{5-6}
% \multicolumn{2}{c|}{} &  &  & \textbf{M4} & \textbf{FR} \\
% \hline
% \multirow{3}{*}{M3}
% & Yearly    & 22.14 & 22.50 & 16.74 & 16.97 \\
% % \cline{2-6}
% & Quarterly & 17.67 & 17.24 & 10.09 & 10.62 \\
% % \cline{2-6}
% & Monthly   & 21.05 & 21.77 & 16.04 & 16.90 \\
% \hline
% \multirow{3}{*}{M4}
% & Yearly    & 36.48 & 34.91 & --    & 14.84 \\
% % \cline{2-6}
% & Quarterly & 34.48 & 34.93 & --    & 12.02 \\
% % \cline{2-6}
% & Monthly   & 31.44 & 32.86 & --    & 15.80 \\
% \hline
% \end{tabular}
% \end{table}

% \begin{table}
%   \caption{Short-horizon forecasting results measured using sMAPE. We compare Raw-$k$NN and PCA-$k$NN transfer results with MOMENT results from Table 3 of \cite{goswami2024moment} under the same source-to-target evaluation setting.}
%   \label{tab:short_horizon_results_transfer}
%   \centering
%   \small
%   \begin{tabular}{@{}llcccc@{}}
%     \toprule
%     \multicolumn{2}{c}{Setting} &
%     Raw-$k$NN &
%     PCA-$k$NN &
%     \multicolumn{2}{c}{MOMENT$_{\mathrm{LP}}$} \\
%     \cmidrule(lr){1-2}
%     % \cmidrule(lr){3-4}
%     \cmidrule(lr){5-6}
%     Dataset & Frequency & & & M4 & FR \\
%     \midrule
%     \multirow{3}{*}{M3}
%       & Yearly    & 22.14 & 22.50 & 16.74 & 16.97 \\
%       & Quarterly & 17.67 & 17.24 & 10.09 & 10.62 \\
%       & Monthly   & 21.05 & 21.77 & 16.04 & 16.90 \\
%     \midrule
%     \multirow{3}{*}{M4}
%       & Yearly    & 36.48 & 34.91 & --    & 14.84 \\
%       & Quarterly & 34.48 & 34.93 & --    & 12.02 \\
%       & Monthly   & 31.44 & 32.86 & --    & 15.80 \\
%     \bottomrule
%   \end{tabular}
% \end{table}

\begin{table}
  \caption{Short-horizon forecasting results measured using sMAPE. We compare Raw-$k$NN and PCA-$k$NN transfer results with MOMENT results from Table 3 of \cite{goswami2024moment} under the same source-to-target evaluation setting.}
  \label{tab:short_horizon_results_transfer}
  \centering
  \small
  \begin{tabular}{@{}llcccc@{}}
    \toprule
    \multicolumn{2}{c}{Setting} &
    \multirow{2}{*}{Raw-$k$NN} &
    \multirow{2}{*}{PCA-$k$NN} &
    \multicolumn{2}{c}{MOMENT$_{\mathrm{LP}}$} \\
    \cmidrule(lr){1-2}
    \cmidrule(lr){5-6}
    Dataset & Frequency & & & M4 & FR \\
    \midrule
    \multirow{3}{*}{M3}
      & Yearly    & 22.14 & 22.50 & 16.74 & 16.97 \\
      & Quarterly & 17.67 & 17.24 & 10.09 & 10.62 \\
      & Monthly   & 21.05 & 21.77 & 16.04 & 16.90 \\
    \midrule
    \multirow{3}{*}{M4}
      & Yearly    & 36.48 & 34.91 & --    & 14.84 \\
      & Quarterly & 34.48 & 34.93 & --    & 12.02 \\
      & Monthly   & 31.44 & 32.86 & --    & 15.80 \\
    \bottomrule
  \end{tabular}
\end{table}


% \begin{table}[H]
% \centering
% \caption{Short-horizon forecasting results measured using sMAPE under the within-group evaluation setting. We report our Raw-$k$NN and PCA-$k$NN results together with MOMENT results from \citep{goswami2024moment} in Table 3.}
% \label{tab:short_horizon_results_within_group }
% \small
% \begin{tabular}{l|l|c|c|c|c}
% \hline
% \multicolumn{2}{c|}{\multirow{2}{*}{\textbf{Dataset}}} & \multirow{2}{*}{\textbf{Raw-$k$NN}} & \multirow{2}{*}{\textbf{PCA-$k$NN}} & \multicolumn{2}{c}{\textbf{MOMENT$_{\mathrm{LP}}$}} \\
% \cline{5-6}
% \multicolumn{2}{c|}{} &  &  & \textbf{M4} & \textbf{FR} \\
% \hline
% \multirow{3}{*}{M3}
% & Yearly    & 20.56 & 21.32 & 16.74 & 16.97 \\
% % \cline{2-6}
% & Quarterly & 14.33 & 14.38 & 10.09 & 10.62 \\
% % \cline{2-6}
% & Monthly   & 15.78 & 16.17 & 16.04 & 16.90 \\
% \hline
% \multirow{3}{*}{M4}
% & Yearly    & 16.65 & 16.69 & --    & 14.84 \\
% % \cline{2-6}
% & Quarterly & 14.51 & 14.84 & --    & 12.02 \\
% % \cline{2-6}
% & Monthly   & 15.37 & 15.53 & --    & 15.80 \\
% \hline
% \end{tabular}
% \end{table}

% \begin{table}
%   \caption{Short-horizon forecasting results measured using sMAPE under the within-group evaluation setting. We compare Raw-$k$NN and PCA-$k$NN with MOMENT results from Table 3 of \cite{goswami2024moment}.}
%   \label{tab:short_horizon_results_within_group }
%   \centering
%   \small
%   \begin{tabular}{@{}llcccc@{}}
%     \toprule
%     \multicolumn{2}{c}{Target series} &
%     Raw-$k$NN &
%     PCA-$k$NN &
%     \multicolumn{2}{c}{MOMENT$_{\mathrm{LP}}$} \\
%     \cmidrule(lr){1-2}
%     \cmidrule(lr){5-6}
%     Dataset & Frequency & & & M4 & FR \\
%     \midrule
%     \multirow{3}{*}{M3}
%       & Yearly    & 20.56 & 21.32 & 16.74 & 16.97 \\
%       & Quarterly & 14.33 & 14.38 & 10.09 & 10.62 \\
%       & Monthly   & 15.78 & 16.17 & 16.04 & 16.90 \\
%     \midrule
%     \multirow{3}{*}{M4}
%       & Yearly    & 16.65 & 16.69 & --    & 14.84 \\
%       & Quarterly & 14.51 & 14.84 & --    & 12.02 \\
%       & Monthly   & 15.37 & 15.53 & --    & 15.80 \\
%     \bottomrule
%   \end{tabular}
% \end{table}


\begin{table}
  \caption{Short-horizon forecasting results measured using sMAPE under the within-group evaluation setting. We compare Raw-$k$NN and PCA-$k$NN with MOMENT results from Table 3 of \cite{goswami2024moment}.}
  \label{tab:short_horizon_results_within_group }
  \centering
  \small
  \begin{tabular}{@{}llcccc@{}}
    \toprule
    \multicolumn{2}{c}{Target series} &
    \multirow{2}{*}{Raw-$k$NN} &
    \multirow{2}{*}{PCA-$k$NN} &
    \multicolumn{2}{c}{MOMENT$_{\mathrm{LP}}$} \\
    \cmidrule(lr){1-2}
    \cmidrule(lr){5-6}
    Dataset & Frequency & & & M4 & FR \\
    \midrule
    \multirow{3}{*}{M3}
      & Yearly    & 20.56 & 21.32 & 16.74 & 16.97 \\
      & Quarterly & 14.33 & 14.38 & 10.09 & 10.62 \\
      & Monthly   & 15.78 & 16.17 & 16.04 & 16.90 \\
    \midrule
    \multirow{3}{*}{M4}
      & Yearly    & 16.65 & 16.69 & --    & 14.84 \\
      & Quarterly & 14.51 & 14.84 & --    & 12.02 \\
      & Monthly   & 15.37 & 15.53 & --    & 15.80 \\
    \bottomrule
  \end{tabular}
\end{table}


% \subsubsection{Imputation Results}

% \begin{table}[H]
% \centering
% \caption{Imputation results. We report test MSE and MAE for Raw-$k$NN, PCA-$k$NN, and MOMENT. MOMENT results are taken from \citep{goswami2024moment} in the Table 6.}
% \label{tab:imputation_results}
% \small
% \begin{tabular}{C{2.2cm}|cc|cc|cc|cc}
% \toprule
% \multirow{2}{*}{Dataset}  & \multicolumn{2}{c|}{Raw-$k$NN} & \multicolumn{2}{c|}{PCA-$k$NN} & \multicolumn{2}{c|}{MOMENT$_0$} & \multicolumn{2}{c}{MOMENT$_{\mathrm{LP}}$} \\
% & MSE & MAE & MSE & MAE & MSE & MAE & MSE & MAE \\
% \midrule
% Weather     & 0.356 & 0.341 & 0.338 & 0.341 & 0.082 & 0.130 & 0.035 & 0.075 \\
% ETTh1       & 0.871 & 0.700 & 0.878 & 0.733 & 0.402 & 0.403 & 0.139 & 0.234 \\
% ETTh2       & 1.571 & 0.891 & 1.547 & 0.884 & 0.125 & 0.238 & 0.061 & 0.159 \\
% ETTm1       & 0.628 & 0.571 & 0.617 & 0.612 & 0.202 & 0.288 & 0.074 & 0.168 \\
% ETTm2       & 1.060 & 0.730 & 0.908 & 0.682 & 0.078 & 0.184 & 0.031 & 0.108 \\
% Electricity & 0.801 & 0.684 & 0.781 & 0.676 & 0.250 & 0.371 & 0.094 & 0.211 \\
% \bottomrule
% \end{tabular}
% \end{table}

\subsection{Dataset availability}\label{Appendix: Dataset_availability}
We use two types of datasets in this paper: synthetic datasets and real-world time-series benchmark datasets. The synthetic datasets in Experiments~\ref{exp1:Moment convergence rate} and~\ref{exp2: Entropy estimation} are generated from the data-generating processes described in Appendices~\ref{app: Experiment 1 implementation details} and~\ref{app: Experiment 2 implementation details}, respectively, using the global random seed 20260504. The real-world datasets in Experiment~\ref{exp3: Real-data_evaluation} are taken from the Time Series PILE~\citep{goswami2024moment}, which includes datasets from the Informer long-horizon forecasting benchmark, the Monash time series forecasting archive, and the UCR classification archive. These benchmark datasets are publicly available from their original sources.

\subsection{Code availability}\label{Appendix: Code_availability}
Code for synthetic data generation, preprocessing, hyperparameter tuning, evaluation, and figure reproduction will be made publicly available upon publication. Experiments~\ref{exp1:Moment convergence rate} and~\ref{exp2: Entropy estimation} were run entirely on CPU on a cluster node with two AMD EPYC 9754 128-Core Processors, comprising 256 physical cores and 256 hardware threads. Reproducing these two experiments takes approximately 5 hours. Experiment~\ref{exp3: Real-data_evaluation} was run entirely on CPU on a Linux workstation with an AMD EPYC 9554 64-Core Processor, comprising 64 physical cores and 128 hardware threads, and takes approximately 4 hours to reproduce. No GPU acceleration was used.


\end{document}
